%===============================================================
%  PART IV — Generative Deep Learning
%===============================================================
\part{Generative Deep Learning}
\label{part:generative}

Questa parte introduce i modelli generativi profondi: dall'autoencoder regolarizzato ai modelli moderni di trasporto di probabilità (flow matching). Il filo conduttore è duplice. Da un lato, una tassonomia comune (figura nella slide GDL23) divide i modelli in \emph{esplicitamente densi} (VAE, normalizing flows, autoregressivi, diffusion via ELBO) e \emph{implicitamente densi} (GAN, modelli che imparano un processo di campionamento). Dall'altro, la famiglia di modelli vista dal punto di vista della \emph{trasformazione di rumore in dati}: VAE $\to$ GAN $\to$ Flows $\to$ Diffusion $\to$ Flow Matching costituiscono uno spettro in cui ogni modello rilassa o specializza il precedente.

%===============================================================
\chapter{Autoencoders e Manifold Learning}
\label{ch:ae}

\begin{risorsevideo}{GDL23 — Autoencoder e manifold hypothesis}
\videolezione{della lezione GDL23}

\medskip
\video{https://www.youtube.com/watch?v=hZ4a4NgM3u0}{Autoencoders — Deep Learning Animated}{Deepia, ~20m}\\
Divulgativo e visivamente ottimo per l'intuizione geometrica (bottleneck, manifold), ma non copre il taglio ``regularized AE'' delle slide: match dichiarato parziale.

\medskip
\video{https://www.youtube.com/watch?v=2859tNY-G5E}{Regularized Autoencoders: Denoising, Contractive, Variational}{Charu Aggarwal, ~40m}\\
Questo invece copre esattamente sparse, denoising e contractive con le rispettive regolarizzazioni: è il complemento del primo.
\end{risorsevideo}

\section{Motivazione e posizionamento}

L'obiettivo della modellazione generativa è apprendere, da un dataset $\mathcal{D} = \{\vect{x}_1, \dots, \vect{x}_N\}$ campionato da una distribuzione sconosciuta $P(\vect{x})$, un modello $P_\theta(\vect{x}) \approx P(\vect{x})$ oppure un processo di campionamento che produca esemplari indistinguibili da quelli in $\mathcal{D}$. Esistono due famiglie:
\begin{description}
  \item[Modelli espliciti] apprendono $P_\theta(\vect{x})$ in forma chiusa o tramite un \emph{lower bound} (VAE, flows, diffusion, autoregressivi).
  \item[Modelli impliciti] apprendono solo una procedura di campionamento $\vect{x} \sim P_\theta$ senza esprimere la densità (GAN, generative stochastic networks).
\end{description}

L'autoencoder (AE) è il punto di partenza: storicamente nato come modello \emph{non} generativo per riduzione di dimensionalità, viene rimodulato per generare attraverso interpretazioni probabilistiche (DAE come stimatore della score, VAE come modello a variabile latente differenziabile).

\section{Autoencoder neurale}

\begin{definition}[Autoencoder]
Un autoencoder è una coppia di funzioni $(f_\phi, g_\theta)$ con $f_\phi: \mathcal{X} \to \mathcal{Z}$ (encoder) e $g_\theta: \mathcal{Z} \to \mathcal{X}$ (decoder), addestrata a minimizzare un errore di ricostruzione
\[
\mathcal{L}_{\text{AE}}(\theta,\phi) = \frac{1}{N}\sum_{i=1}^N L\big(\vect{x}_i, g_\theta(f_\phi(\vect{x}_i))\big),
\]
tipicamente con $L = \|\cdot\|_2^2$. La rappresentazione $\vect{z} = f_\phi(\vect{x})$ vive nel \emph{bottleneck}, con dim$(\mathcal{Z}) \ll$ dim$(\mathcal{X})$, oppure (in AE sovracompleti) con un vincolo di regolarizzazione che impedisce l'apprendimento dell'identità.
\end{definition}

\begin{remark}[Manifold hypothesis]
Si assume che i dati osservati giacciano (approssimativamente) su una sottovarietà $\mathcal{M} \subset \mathcal{X}$ di dimensione molto inferiore a $\dim(\mathcal{X})$. L'AE impara una parametrizzazione di $\mathcal{M}$ via $f_\phi$ (immersione) e $g_\theta$ (ricostruzione). La regolarizzazione spinge $g_\theta \circ f_\phi$ a essere sensibile solo alle variazioni \emph{tangenti} alla varietà.
\end{remark}

\begin{intuizione}{il bottleneck non è l'unico modo di impedire la copia}
Il rischio strutturale di un autoencoder è banale: se encoder e decoder sono abbastanza espressivi, la soluzione ottima è l'identità. Ricostruzione perfetta, rappresentazione inutile.

Il bottleneck (latente più piccolo dell'input) è il rimedio ovvio: se lo spazio non basta, la rete deve buttare via qualcosa, e conviene che butti via il rumore invece del segnale. Ma è un rimedio grezzo, perché la capacità va scelta a mano e non c'è garanzia che ciò che sopravvive sia semanticamente interessante.

Le varianti regolarizzate sostituiscono il vincolo di \emph{dimensione} con un vincolo sul \emph{comportamento}, e questo permette di usare anche latenti sovracompleti:
\begin{itemize}[leftmargin=1.4em]
  \item \textbf{Sparse}: penalizza $\|\vect{h}\|_1$. Il latente può essere grande, ma solo poche unità si accendono per input — capacità limitata per campione, non in assoluto.
  \item \textbf{Denoising}: corrompe l'input e chiede di ricostruire l'originale. L'identità non è più una soluzione, perché copierebbe anche il rumore. La rete è costretta a imparare che cosa \emph{non} appartiene ai dati.
  \item \textbf{Contractive}: penalizza $\|\partial \vect{h}/\partial \vect{x}\|_F^2$. Chiede esplicitamente che la rappresentazione sia insensibile a piccole perturbazioni, tranne che nelle direzioni tangenti alla varietà dei dati.
\end{itemize}

\textbf{Il filo che porta al resto della Parte IV}: DAE e CAE arrivano alla stessa conclusione da direzioni diverse, cioè che il modello deve imparare come i dati si comportano \emph{attorno} alla varietà, non solo sopra. Il DAE in particolare, si dimostrerà, stima implicitamente $\nabla_{\vect{x}} \log p(\vect{x})$ — la \emph{score}. È il motivo per cui il capitolo su diffusion e score matching non è un argomento nuovo, ma il seguito di questo.
\end{intuizione}

\section{Varianti regolarizzate}

\subsection{Sparse autoencoder (SAE)}
Aggiunge una penalità sull'attivazione latente
\[
\mathcal{L}_{\text{SAE}} = L(\vect{x}, \tilde{\vect{x}}) + \lambda \, \Omega(\vect{z}), \qquad \Omega(\vect{z}) = \|\vect{z}\|_1 = \sum_j |z_j|.
\]
\begin{remark}[Interpretazione MAP]
$\Omega(\vect{z}) = \lambda \|\vect{z}\|_1$ corrisponde a un prior di Laplace $P(\vect{z}) \propto \exp(-\lambda \|\vect{z}\|_1/2)$. Allenare con regolarizzazione equivale (a meno di costanti) a $\max_\theta \log P(\vect{x}|\vect{z}) + \log P(\vect{z})$, ossia inferenza MAP.
\end{remark}

\subsection{Denoising autoencoder (DAE)}
Si corrompe l'input con un processo di rumore $C(\hat{\vect{x}} | \vect{x})$, tipicamente $\hat{\vect{x}} = \vect{x} + \boldsymbol{\epsilon}$ con $\boldsymbol{\epsilon} \sim \mathcal{N}(0, \sigma^2 \mat{I})$, e si minimizza
\[
\mathcal{L}_{\text{DAE}} = \mathbb{E}_{\vect{x},\hat{\vect{x}}}\, \| \vect{x} - g_\theta(f_\phi(\hat{\vect{x}})) \|_2^2 .
\]

\begin{theorem}[DAE come stima della score, Vincent 2011]
Sotto rumore gaussiano $\hat{\vect{x}} = \vect{x} + \boldsymbol{\epsilon}$ con $\boldsymbol{\epsilon} \sim \mathcal{N}(0, \sigma^2 \mat{I})$ e per $\sigma$ piccolo, il minimo della loss DAE è il \emph{conditional mean estimator} $g(f(\hat{\vect{x}})) = \mathbb{E}[\vect{x} \mid \hat{\vect{x}}]$. Applicando la formula di Tweedie
\[
\mathbb{E}[\vect{x} \mid \hat{\vect{x}}] = \hat{\vect{x}} + \sigma^2 \nabla_{\hat{\vect{x}}} \log p(\hat{\vect{x}}),
\]
si ottiene
\[
\boxed{\;\nabla_{\hat{\vect{x}}} \log p(\hat{\vect{x}}) = \frac{g_\theta(f_\phi(\hat{\vect{x}})) - \hat{\vect{x}}}{\sigma^2}.\;}
\]
Il campo vettoriale $g(f(\hat{\vect{x}})) - \hat{\vect{x}}$ approssima il gradiente della log-densità (la \emph{score function}).
\end{theorem}

\begin{remark}[Ponte con diffusion]
Questa identità è il seme del legame fra DAE, score matching e modelli di diffusione (Cap.~\ref{ch:diffusion} e \ref{ch:flowmatching}): un DAE già impara la score; un diffusion model è un DAE gerarchico a livelli multipli di $\sigma$.
\end{remark}

\subsection{Contractive autoencoder (CAE)}
Penalizza la sensibilità dell'encoder
\[
\mathcal{L}_{\text{CAE}} = L(\vect{x}, \tilde{\vect{x}}) + \lambda \left\| \frac{\partial f_\phi(\vect{x})}{\partial \vect{x}} \right\|_F^2 .
\]
La rappresentazione $\vect{z}$ varia poco rispetto a perturbazioni infinitesime di $\vect{x}$: si apprende un'immersione locale a varietà.

\section{Deep autoencoders e pretraining}

Un deep AE è una composizione gerarchica di encoder $f_1, \dots, f_L$ e decoder simmetrici $g_L, \dots, g_1$. Storicamente, prima di Adam e batch normalization, l'addestramento end-to-end soffriva del vanishing gradient (Cap.~\ref{ch:propagation}). La soluzione era il \emph{layerwise unsupervised pretraining}:
\begin{enumerate}
  \item Addestra $(W_1)$ come AE su $\vect{x}$;
  \item Congela $W_1$, addestra $(W_2)$ come AE su $\vect{z}_1 = f_1(\vect{x})$;
  \item Continua fino al livello $L$;
  \item (Opzionale) fine-tuning end-to-end di tutta la pila.
\end{enumerate}

\paragraph{Connessione con RBM.}
Riscrivendo lo stack di autoencoder come una struttura a strati di unità stocastiche bidirezionali si ottiene la \emph{Deep Belief Network} (DBN): una pila di RBM (Cap.~\ref{ch:mrf}) addestrate con contrastive divergence. La DBN è (in larga parte) un grafo diretto top-down con uno strato superiore non diretto (l'ultimo RBM). La \emph{Deep Boltzmann Machine} (DBM), invece, è un modello completamente non diretto, e richiede il \emph{double counting trick} (dimezzare i pesi degli strati intermedi durante il pretraining per evitare di contare due volte i contributi $\vect{x} \to \vect{z}_1$).

\section{Applicazioni}

\begin{description}
  \item[Visualizzazione e indexing] proiezione su 2D/3D per esplorazione dati.
  \item[Anomaly detection] punti con errore di ricostruzione $\|\vect{x} - g(f(\vect{x}))\|$ sopra una soglia sono potenziali anomalie.
  \item[Pretraining supervisato] iniziare un classificatore profondo dai pesi di un deep AE.
  \item[Compressione e denoising] usando CNN/U-Net come encoder/decoder.
  \item[Multimodal learning] DBM multimodale (Srivastava \& Salakhutdinov 2014) fonde più modalità (immagini + testo) in uno spazio condiviso.
\end{description}

\begin{tcolorbox}[mybox, title={Quando usare un autoencoder}]
\begin{description}[leftmargin=4em,style=nextline]
  \item[Si vuole una rappresentazione compatta non supervisionata] DAE o CAE: la regolarizzazione fissa la varietà.
  \item[Anomaly detection] usare la distribuzione degli errori di ricostruzione sui normali.
  \item[Generazione di qualità] \textbf{Non} l'AE classico: l'output del bottleneck non è una distribuzione campionabile. Passare a VAE (Cap.~\ref{ch:vae}) o diffusion (Cap.~\ref{ch:diffusion}).
  \item[Pretraining moderno] generalmente sostituito da self-supervised contrastivo (SimCLR, BYOL) per le immagini, ma DAE rimane competitivo per dati strutturati e biomedici.
\end{description}
\end{tcolorbox}

\begin{tcolorbox}[mybox, title={Connessioni}]
\begin{itemize}
  \item DAE $\leftrightarrow$ score matching: il DAE apprende il campo $\nabla_{\vect{x}} \log p(\vect{x})$ (Tweedie).
  \item AE regolarizzato $\leftrightarrow$ MAP: la regolarizzazione $\Omega(\vect{z})$ implementa un prior sulla latente.
  \item Deep AE $\leftrightarrow$ DBN/DBM: la stessa architettura, ma con unità stocastiche e training contrastivo.
  \item AE $\to$ VAE: aggiungere un encoder stocastico e una loss ELBO trasforma l'AE in un modello latente generativo.
\end{itemize}
\end{tcolorbox}

\begin{tcolorbox}[mywarn, title={Domande d'esame — Autoencoders}]
\begin{enumerate}
  \item Cosa significa che il DAE impara un'approssimazione della score function? Deriva l'identità di Tweedie e mostra come si arriva a $\nabla_{\hat{\vect{x}}} \log p(\hat{\vect{x}}) = (g(f(\hat{\vect{x}})) - \hat{\vect{x}})/\sigma^2$.
  \item Differenza concettuale e formale fra sparse AE, denoising AE e contractive AE. Quando ne preferiresti uno rispetto agli altri?
  \item Spiega il \emph{layerwise pretraining} di un deep AE. Perché era necessario in passato e perché oggi è meno usato?
  \item Differenza fra DBN e DBM: quale è diretto, quale no, come differisce l'inferenza.
  \item Perché un AE classico non è un buon modello generativo, anche se ha un bottleneck? (suggerimento: lo spazio latente non è distribuito secondo una distribuzione campionabile).
  \item Confronto fra Tweedie's formula e reparametrization trick: entrambi permettono di trattare distribuzioni gaussiane in modi che facilitano l'apprendimento — sono lo stesso strumento?
\end{enumerate}
\end{tcolorbox}

%===============================================================
\chapter{Variational Autoencoders}

\begin{risorsevideo}{GDL24 — Variational Autoencoders}
\videolezione{della lezione GDL24}

\medskip
\video{https://www.youtube.com/watch?v=NlIqjtbjjRE}{CS294-158 SP24 — L4: Latent Variable Models and Variational AutoEncoders}{Pieter Abbeel, UC Berkeley, ~2h}\\
Deriva l'ELBO senza saltare passaggi e spiega il reparametrization trick come soluzione a un problema preciso di stima del gradiente. Lungo, ma è la trattazione più completa disponibile.

\medskip
\video{https://www.youtube.com/watch?v=DaqNNLidswA}{Variational Inference: Foundations and Innovations}{David Blei, ~1h30m}\\
Da rivedere se l'ELBO della Parte~II non è ancora automatico: il VAE è quella stessa ELBO con $q$ parametrizzata da una rete.
\end{risorsevideo}
\label{ch:vae}

\section{Dalla DAE al modello a variabile latente}

L'osservazione cruciale (Cap.~\ref{ch:ae}, formula di Tweedie) è che il DAE stima $\nabla_{\vect{x}} \log p(\vect{x})$ solo nell'intorno dei training point. Vogliamo invece un modello esplicito di $p_\theta(\vect{x})$, ben definito su tutto $\mathcal{X}$.

L'idea è un modello a variabile latente continua:
\begin{equation}\label{eq:lvm}
p_\theta(\vect{x}) = \int p_\theta(\vect{x}\mid \vect{z}) \, p(\vect{z})\, d\vect{z},
\end{equation}
con $p(\vect{z}) = \mathcal{N}(0, \mat{I})$ prior semplice e $p_\theta(\vect{x}|\vect{z})$ rete neurale (decoder). L'integrale è intrattabile in dimensione alta: serve un'approssimazione variazionale.

\section{ELBO}

\begin{theorem}[Evidence Lower BOund]
Per qualunque distribuzione $q_\phi(\vect{z}|\vect{x})$ vale
\[
\log p_\theta(\vect{x}) = \underbrace{\mathbb{E}_{q_\phi(\vect{z}|\vect{x})}\!\left[\log \frac{p_\theta(\vect{x},\vect{z})}{q_\phi(\vect{z}|\vect{x})}\right]}_{\mathcal{L}(\vect{x};\theta,\phi)} + \KL{q_\phi(\vect{z}|\vect{x})}{p_\theta(\vect{z}|\vect{x})}.
\]
\end{theorem}
\begin{proof}
$\log p(\vect{x}) = \log \int p(\vect{x},\vect{z})\,d\vect{z} = \log \mathbb{E}_q\!\left[\frac{p(\vect{x},\vect{z})}{q(\vect{z}|\vect{x})}\right]$. La disuguaglianza di Jensen su $\log$ (concavo) dà $\log p(\vect{x}) \geq \mathbb{E}_q\!\left[\log \frac{p(\vect{x},\vect{z})}{q(\vect{z}|\vect{x})}\right]$. L'identità con il gap-KL si ottiene scrivendo $p(\vect{x},\vect{z}) = p(\vect{z}|\vect{x}) p(\vect{x})$ e separando i termini.
\end{proof}

Poiché $\KL \geq 0$:
\begin{equation}
\log p_\theta(\vect{x}) \geq \mathcal{L}(\vect{x};\theta,\phi).
\end{equation}
Massimizzare $\mathcal{L}$ rispetto a $(\theta, \phi)$ \emph{contemporaneamente} (i) aumenta la verosimiglianza del modello, (ii) avvicina $q_\phi(\vect{z}|\vect{x})$ alla vera posteriori intrattabile $p_\theta(\vect{z}|\vect{x})$.

Una riscrittura utile:
\begin{equation}\label{eq:elbo-vae}
\boxed{\;\mathcal{L}(\vect{x};\theta,\phi) = \underbrace{\mathbb{E}_{q_\phi(\vect{z}|\vect{x})}\!\left[\log p_\theta(\vect{x}|\vect{z})\right]}_{\text{ricostruzione}} - \underbrace{\KL{q_\phi(\vect{z}|\vect{x})}{p(\vect{z})}}_{\text{regolarizzazione}}.\;}
\end{equation}

\section{Architettura VAE}

\begin{description}
  \item[Encoder] $q_\phi(\vect{z}|\vect{x}) = \mathcal{N}(\boldsymbol{\mu}_\phi(\vect{x}), \mathrm{diag}(\boldsymbol{\sigma}^2_\phi(\vect{x})))$: una rete neurale produce $\boldsymbol{\mu}, \boldsymbol{\sigma}$.
  \item[Decoder] $p_\theta(\vect{x}|\vect{z})$ è tipicamente $\mathcal{N}(g_\theta(\vect{z}), \sigma_x^2 \mat{I})$ (output continuo) o Bernoulli (binari) o categoriale (multinomiali).
  \item[Prior] $p(\vect{z}) = \mathcal{N}(0,\mat{I})$.
\end{description}

\subsection{Reparametrization trick}

Il problema: $\mathbb{E}_{q_\phi}[\cdot]$ richiede campionare $\vect{z}$ da $q_\phi$ che dipende da $\phi$. Il sampling non è differenziabile. La soluzione:
\begin{equation}
\vect{z} = \boldsymbol{\mu}_\phi(\vect{x}) + \boldsymbol{\sigma}_\phi(\vect{x}) \odot \boldsymbol{\epsilon}, \qquad \boldsymbol{\epsilon} \sim \mathcal{N}(0,\mat{I}).
\end{equation}
Ora $\boldsymbol{\epsilon}$ non dipende dai parametri, e $\vect{z}$ è una funzione differenziabile di $(\boldsymbol{\mu}, \boldsymbol{\sigma}, \boldsymbol{\epsilon})$. Quindi
\[
\nabla_\phi \mathbb{E}_{q_\phi}[f(\vect{z})] = \mathbb{E}_{\boldsymbol{\epsilon}}\!\left[\nabla_\phi f\big(\boldsymbol{\mu}_\phi(\vect{x}) + \boldsymbol{\sigma}_\phi(\vect{x}) \odot \boldsymbol{\epsilon}\big)\right],
\]
che si stima con Monte Carlo (in pratica un solo $\boldsymbol{\epsilon}$ per esempio).

\begin{intuizione}{il reparametrization trick risolve un problema di gradiente, non di campionamento}
Campionare $\vect{z} \sim q_\phi(\vect{z} \mid \vect{x})$ non è difficile: è una gaussiana, si sa fare. Il problema è un altro: l'operazione di campionamento è \emph{stocastica e non differenziabile rispetto a $\phi$}. Il gradiente che dovrebbe tornare all'encoder trova un nodo casuale sulla strada e si ferma.

Esisterebbe una via alternativa, lo stimatore REINFORCE (score function estimator), che è non distorto ma ha varianza altissima — in pratica inutilizzabile per addestrare un encoder.

Il trucco della riparametrizzazione sposta la casualità fuori dal cammino del gradiente. Invece di campionare $\vect{z}$ da una distribuzione che dipende da $\phi$, si campiona $\boldsymbol{\epsilon} \sim \mathcal{N}(0, \mat{I})$ — che non dipende da niente — e si costruisce
\[
\vect{z} = \boldsymbol{\mu}_\phi(\vect{x}) + \boldsymbol{\sigma}_\phi(\vect{x}) \odot \boldsymbol{\epsilon}.
\]
Ora $\vect{z}$ è una funzione \emph{deterministica e differenziabile} di $\phi$ e di un rumore esterno. Il grafo computazionale è continuo da capo a fondo e la backpropagation passa normalmente.

\textbf{Il modo di dirlo all'orale}: la casualità non è stata eliminata, è stata spostata da dentro il grafo a un ingresso del grafo. La differenza in varianza dello stimatore è di ordini di grandezza, ed è ciò che rende il VAE addestrabile.

\textbf{Limite implicito}: il trucco funziona perché la gaussiana ammette una riparametrizzazione location-scale. Per latenti discreti non esiste, ed è la ragione per cui servono Gumbel-Softmax, straight-through estimator o VQ-VAE.
\end{intuizione}

\subsection{KL in forma chiusa}

Per $q = \mathcal{N}(\boldsymbol{\mu}, \mathrm{diag}(\boldsymbol{\sigma}^2))$ e $p = \mathcal{N}(0, \mat{I})$:
\begin{equation}\label{eq:kl-gauss}
\boxed{\;\KL{q}{p} = \frac{1}{2}\sum_{j=1}^d\big(\mu_j^2 + \sigma_j^2 - 1 - \log \sigma_j^2\big).\;}
\end{equation}
\begin{proof}[Sketch]
$\KL{\mathcal{N}(\mu_1,\sigma_1^2)}{\mathcal{N}(\mu_2,\sigma_2^2)} = \log\frac{\sigma_2}{\sigma_1} + \frac{\sigma_1^2 + (\mu_1-\mu_2)^2}{2\sigma_2^2} - \frac{1}{2}$. Per gaussiane indipendenti la KL è somma componente per componente. Specializzando con $\mu_2=0, \sigma_2=1$ si ottiene~\eqref{eq:kl-gauss}.
\end{proof}

\subsection{Loss finale}
Combinando~\eqref{eq:elbo-vae},~\eqref{eq:kl-gauss} e la reparametrizzazione:
\[
-\mathcal{L} \approx \frac{1}{2\sigma_x^2}\| \vect{x} - g_\theta(\vect{z}) \|^2 + \frac{1}{2}\sum_j (\mu_j^2 + \sigma_j^2 - 1 - \log \sigma_j^2),
\]
con $\vect{z} = \boldsymbol{\mu}_\phi(\vect{x}) + \boldsymbol{\sigma}_\phi(\vect{x}) \odot \boldsymbol{\epsilon}$.

\begin{esempiosvolto}{i due termini dell'ELBO, con numeri}
La loss del VAE per un campione è
\[
-\mathcal{L} = \underbrace{-\mathbb{E}_{q}\big[\log p_\theta(\vect{x} \mid \vect{z})\big]}_{\text{ricostruzione}} + \underbrace{\KL{q_\phi(\vect{z}|\vect{x})}{p(\vect{z})}}_{\text{regolarizzazione}}.
\]

\medskip
\textbf{Il termine KL, in forma chiusa.} Per $q = \mathcal{N}(\mu, \sigma^2)$ e prior $p = \mathcal{N}(0,1)$ in una dimensione:
\[
\KL{q}{p} = \tfrac{1}{2}\big(\mu^2 + \sigma^2 - 1 - \log \sigma^2\big).
\]
Con $\mu = 0.5$ e $\sigma = 0.8$:
\[
\tfrac{1}{2}\big(0.25 + 0.64 - 1 + 0.4463\big) = \tfrac{1}{2}(0.3363) = \mathbf{0.168}.
\]
Si verifichino i due casi limite: con $\mu = 0$, $\sigma = 1$ la KL è esattamente $0$ (il posterior coincide col prior); con $\sigma \to 0$ il termine $-\log \sigma^2$ diverge, quindi il modello è \emph{penalizzato all'infinito} se prova a rendere l'encoder deterministico.

\medskip
\textbf{Che cosa fa quel $-\log \sigma^2$.} È il termine che impedisce al VAE di collassare in un autoencoder normale. Senza rumore nel latente il modello ricostruirebbe meglio, ma lo spazio latente resterebbe pieno di buchi e il campionamento da $p(\vect{z})$ produrrebbe spazzatura. La KL forza le gaussiane posteriori a sovrapporsi, riempiendo lo spazio.

\medskip
\textbf{Il tiro alla fune, e le due patologie.} I due termini si oppongono:
\begin{itemize}[leftmargin=1.4em]
  \item se la KL vince, si ha \emph{posterior collapse}: $q_\phi(\vect{z}|\vect{x}) \to p(\vect{z})$ per ogni $\vect{x}$, la KL va a $0$, il latente non porta più informazione sull'input e il decoder impara a ignorarlo. Si riconosce guardando la KL: se tende a zero mentre la ricostruzione non migliora, è questo.
  \item se vince la ricostruzione, il latente si sparpaglia e il campionamento a priori peggiora.
\end{itemize}
Il $\beta$-VAE non fa che mettere un peso esplicito su questo equilibrio: $\beta > 1$ spinge verso il disentanglement pagando in qualità di ricostruzione, $\beta < 1$ il contrario.

\textbf{Nota sull'immagine sfocata.} Le ricostruzioni del VAE sono tipicamente sfocate perché il termine di ricostruzione, con decoder gaussiano a varianza fissa, è un errore quadratico: di fronte a un'ambiguità, l'MSE preferisce la \emph{media} di tutte le possibilità plausibili a una qualsiasi di esse. La media di più volti nitidi è un volto sfocato. Il GAN del capitolo seguente non ha questo problema proprio perché non ottimizza una verosimiglianza per pixel.
\end{esempiosvolto}

\section{Generazione e CVAE}

\paragraph{Campionamento (test).} Si stacca l'encoder: $\vect{z} \sim \mathcal{N}(0,\mat{I})$, $\tilde{\vect{x}} = g_\theta(\vect{z})$.

\paragraph{Conditional VAE.} Si aggiunge un input $\vect{y}$ (etichetta o condizione) a encoder e decoder:
\[
\mathcal{L}_{\text{CVAE}}(\vect{x},\vect{y}) = \mathbb{E}_{q_\phi(\vect{z}|\vect{x},\vect{y})}[\log p_\theta(\vect{x}|\vect{z},\vect{y})] - \KL{q_\phi(\vect{z}|\vect{x},\vect{y})}{p(\vect{z}|\vect{y})}.
\]

\section{Interpretazione info-teorica}

Riscrivendo l'ELBO:
\[
\mathcal{L} = \mathbb{E}_{\vect{x}\sim\mathcal{D}}\!\left[\mathbb{E}_{q_\phi(\vect{z}|\vect{x})}[\log p_\theta(\vect{x}|\vect{z})] - \KL{q_\phi(\vect{z}|\vect{x})}{p(\vect{z})}\right].
\]
\begin{itemize}
  \item Il primo termine è il numero di bit necessari per ricodificare $\vect{x}$ dato $\vect{z}$ con $p_\theta(\vect{x}|\vect{z})$.
  \item Il secondo è l'\emph{information gain}: quanti bit servono per trasmettere $\vect{z}$ usando $q_\phi(\vect{z}|\vect{x})$ rispetto al prior. È anche un costo di codifica.
\end{itemize}
Il VAE è quindi un codec a tasso minimo (compressione lossy con prior gaussiano).

\section{Patologie e varianti}

\subsection{$\beta$-VAE}
\[
\mathcal{L}_\beta = \mathbb{E}_{q_\phi}[\log p_\theta(\vect{x}|\vect{z})] - \beta\, \KL{q_\phi(\vect{z}|\vect{x})}{p(\vect{z})}, \qquad \beta > 1.
\]
Spinge i fattori latenti a essere più indipendenti (\emph{disentanglement}). Vedi Cap.~\ref{ch:crl}.

\subsection{Posterior collapse}
Se il decoder $p_\theta(\vect{x}|\vect{z})$ è troppo espressivo (es.\ autoregressive PixelCNN), può ricostruire $\vect{x}$ \emph{senza usare} $\vect{z}$. Allora $q_\phi(\vect{z}|\vect{x}) \to p(\vect{z})$, il KL va a zero e la latente diventa inutile. Mitigazioni: KL annealing, free bits, $\beta < 1$ iniziale.

\subsection{Confronto con AE classico e DAE}
\begin{itemize}
  \item L'AE classico non ha vincoli sul prior di $\vect{z}$: campionare random non ha senso.
  \item Il DAE apprende la score \emph{solo} attorno ai training point. Il VAE definisce $p_\theta(\vect{x})$ globalmente.
\end{itemize}

\begin{tcolorbox}[mybox, title={Quando usare un VAE}]
\begin{description}[leftmargin=4em,style=nextline]
  \item[Representation learning con prior strutturato] sì, è il \emph{best in class} per spazi latenti regolari, interpolazioni, manipolazione semantica.
  \item[Generazione fotorealistica] \textbf{no}: campioni tipicamente sfocati. Usare diffusion o GAN.
  \item[Backbone modulare] sì, il VAE è il building block di latent diffusion (Stable Diffusion), CITRIS, e di buona parte dei modelli causali.
  \item[Anomaly detection] sì, l'ELBO è una misura naturale di tipicità.
  \item[Conditional generation] facile via CVAE.
\end{description}
\end{tcolorbox}

\begin{tcolorbox}[mybox, title={Connessioni}]
\begin{itemize}
  \item ELBO $=$ stessa quantità di Variational Inference (Cap.~\ref{ch:variational}); il reparametrization trick rende l'EM differenziabile e parametrizza encoder/decoder con reti neurali.
  \item Il VAE è il padre di diffusion (visto come VAE gerarchico con encoder fisso) e di Causal Representation Learning (iVAE, $\beta$-VAE, CITRIS).
  \item DAE $\to$ VAE: passare da score locale a densità globale.
  \item Reparametrization trick $\to$ flow matching: in entrambi si differenzia attraverso un campionamento.
\end{itemize}
\end{tcolorbox}

\begin{tcolorbox}[mywarn, title={Domande d'esame — VAE}]
\begin{enumerate}
  \item Deriva l'ELBO partendo da $\log p_\theta(\vect{x}) = \log \int p(\vect{x},\vect{z})\,d\vect{z}$ usando la disuguaglianza di Jensen. Mostra esplicitamente come compaiono i termini di ricostruzione e KL.
  \item Cos'è il reparametrization trick e perché è necessario? Cosa succederebbe se rimuovi questo trucco e provi a backpropagare attraverso il campionamento?
  \item Deriva $\KL{\mathcal{N}(\boldsymbol{\mu}, \mathrm{diag}(\boldsymbol{\sigma}^2))}{\mathcal{N}(0,\mat{I})}$ in forma chiusa.
  \item Cos'è il \emph{posterior collapse}? In quali condizioni si verifica e come si mitiga?
  \item Confronta VAE con DAE: entrambi usano rumore gaussiano, dov'è la differenza?
  \item Perché i campioni VAE sono tipicamente sfocati? È un problema dell'ELBO o della parametrizzazione gaussiana del decoder?
  \item Come modificheresti l'ELBO per condizionare la generazione su un'etichetta $\vect{y}$? Scrivi la loss CVAE.
\end{enumerate}
\end{tcolorbox}

%===============================================================
\chapter{Generative Adversarial Networks}
\label{ch:gan}

\begin{risorsevideo}{GDL25 — Generative Adversarial Networks}
\videolezione{della lezione GDL25}

\medskip
\video{https://www.youtube.com/watch?v=lFAHPJS2HHc}{CS294-158 SP24 — L5: GANs}{Pieter Abbeel, UC Berkeley, ~2h}\\
Deriva $D^*$ e la connessione con la Jensen-Shannon esattamente come le slide, e prosegue con WGAN, mode collapse e le varianti condizionali. È il riferimento completo del capitolo.
\end{risorsevideo}

\section{Idea: imparare a campionare}

A differenza del VAE che apprende esplicitamente $p_\theta(\vect{x})$ via ELBO, una GAN \emph{non} apprende mai la densità. Invece, apprende un \emph{generatore} $G_{\theta_G}: \mathcal{Z} \to \mathcal{X}$ (rete deterministica) tale che $G(\vect{z})$ con $\vect{z} \sim p(\vect{z}) = \mathcal{N}(0,\mat{I})$ produca campioni \emph{indistinguibili} dai dati reali agli occhi di un \emph{discriminatore} $D_{\theta_D}: \mathcal{X} \to [0,1]$.

\section{Formulazione minimax}

\begin{equation}\label{eq:gan-minmax}
\boxed{\;\min_{\theta_G} \max_{\theta_D}\; V(D, G) = \mathbb{E}_{\vect{x}\sim p_{\text{data}}}[\log D(\vect{x})] + \mathbb{E}_{\vect{z}\sim p(\vect{z})}[\log(1 - D(G(\vect{z})))].\;}
\end{equation}

$D$ massimizza la probabilità di etichettare correttamente reali e falsi; $G$ minimizza la probabilità che $D$ smascheri i suoi campioni.

\subsection{Discriminatore ottimale}

\begin{theorem}[Goodfellow et al.\ 2014]
Per $G$ fisso, il discriminatore ottimale di~\eqref{eq:gan-minmax} è
\[
\boxed{\;D^*(\vect{x}) = \frac{p_{\text{data}}(\vect{x})}{p_{\text{data}}(\vect{x}) + p_G(\vect{x})},\;}
\]
dove $p_G$ è la densità implicita indotta da $G$.
\end{theorem}
\begin{proof}
Per ogni $\vect{x}$ fissato, il funzionale è $f(d) = p_{\text{data}}(\vect{x})\log d + p_G(\vect{x})\log(1-d)$. Derivando rispetto a $d$ e azzerando: $p_{\text{data}}/d - p_G/(1-d) = 0 \Rightarrow d = p_{\text{data}}/(p_{\text{data}} + p_G)$. La seconda derivata è negativa, dunque è un massimo.
\end{proof}

\subsection{Connessione con Jensen-Shannon}

Sostituendo $D^*$ in $V$:
\begin{align*}
V(D^*, G) &= \mathbb{E}_{p_{\text{data}}}\!\left[\log\frac{p_{\text{data}}}{p_{\text{data}}+p_G}\right] + \mathbb{E}_{p_G}\!\left[\log\frac{p_G}{p_{\text{data}}+p_G}\right]\\
&= -\log 4 + \KL{p_{\text{data}}}{\tfrac{p_{\text{data}}+p_G}{2}} + \KL{p_G}{\tfrac{p_{\text{data}}+p_G}{2}}\\
&= -\log 4 + 2 \cdot \mathrm{JS}(p_{\text{data}} \| p_G).
\end{align*}
Quindi il problema GAN minimizza la divergenza di Jensen-Shannon fra $p_{\text{data}}$ e $p_G$. L'ottimo globale è $p_G = p_{\text{data}}$, raggiunto con $V = -\log 4$ e $D^* \equiv 1/2$ (indistinguibile).

\begin{intuizione}{la GAN non minimizza una loss, cerca un equilibrio}
La differenza radicale con tutto il resto della Parte IV: VAE, flow e diffusion minimizzano una funzione obiettivo. Una GAN no. Il problema
\[
\min_G \max_D V(D,G)
\]
è la ricerca di un \emph{punto di sella}, e questo cambia tutto ciò che si può dire sul training.

Non esiste una quantità che scende monotonamente e che si possa guardare per capire se le cose vanno bene: la loss del generatore può salire mentre i campioni migliorano, e viceversa. Non c'è un criterio di stop naturale. La dinamica può oscillare o entrare in ciclo, perché due gradienti che si contrastano non convergono necessariamente — anche nel caso quadratico più semplice il gradient descent-ascent può orbitare intorno alla soluzione senza raggiungerla.

\textbf{Perché comunque funziona (quando funziona)}: sostituendo $D^*$ nella $V$ si ottiene
\[
V(D^*, G) = -\log 4 + 2\, \mathrm{JSD}(p_{\text{data}} \| p_G),
\]
quindi \emph{all'ottimo del discriminatore} il generatore sta effettivamente minimizzando una divergenza, e il minimo $-\log 4 \approx -1.386$ si raggiunge sse $p_G = p_{\text{data}}$. Il punto sottile è quel ``all'ottimo'': nella pratica $D$ non è mai ottimo, quindi il generatore minimizza un'approssimazione della JSD, e le garanzie teoriche non si applicano.

\textbf{E il vantaggio, che va detto insieme ai difetti}: proprio perché non ottimizza una verosimiglianza per pixel, la GAN non è spinta a mediare fra le alternative plausibili. È il motivo per cui i suoi campioni sono nitidi dove quelli del VAE sono sfocati — e anche il motivo per cui può ignorare interi modi della distribuzione senza pagare pegno, che è il mode collapse.
\end{intuizione}

\section{Problemi pratici}

\subsection{Gradient saturation}
Quando il generatore è scarso, $D(G(\vect{z})) \to 0$ e $\log(1-D(G(\vect{z}))) \to 0$ con gradiente quasi piatto. La fix (Goodfellow 2014) è il \emph{non-saturating loss}: il generatore massimizza $\mathbb{E}_{\vect{z}}[\log D(G(\vect{z}))]$ invece di minimizzare $\mathbb{E}_{\vect{z}}[\log(1 - D(G(\vect{z})))]$. Stesso punto fisso, gradiente meno piatto.

\subsection{Mode collapse}
Il generatore impara a produrre pochi modi della distribuzione che \emph{ingannano} $D$, ignorando il resto. Sintomo: campioni simili tra loro.

\subsection{Instabilità del training}
Il problema minimax è un punto sella, non un minimo. Senza heuristiche (batch normalization, label smoothing, spectral norm, two-timescale rules) il training oscilla.

\begin{esempiosvolto}{gradient saturation: perché si cambia la loss del generatore}
La formulazione originale chiede al generatore di minimizzare $\log(1 - D(G(\vect{z})))$. Vediamo che cosa succede all'inizio del training, quando $D$ smaschera facilmente i campioni falsi e quindi $D(G(\vect{z})) \approx 0.01$.

\medskip
\textbf{Loss saturante.} La derivata rispetto all'uscita del discriminatore è
\[
\left|\frac{d}{dD} \log(1 - D)\right| = \frac{1}{1 - D} = \frac{1}{0.99} \approx 1.01.
\]

\textbf{Loss non saturante.} Si massimizza invece $\log D(G(\vect{z}))$, cioè si minimizza $-\log D$:
\[
\left|\frac{d}{dD}\big(-\log D\big)\right| = \frac{1}{D} = \frac{1}{0.01} = 100.
\]

Un fattore $\approx 99$ nel segnale che arriva al generatore, \emph{nello stesso punto operativo}. E il rapporto peggiora ulteriormente man mano che $D$ migliora: a $D = 10^{-3}$ diventa circa $1000$.

\medskip
\textbf{Perché è esattamente il momento sbagliato per avere gradiente piccolo.} All'inizio del training il generatore produce spazzatura, quindi $D$ vince facilmente: è proprio la fase in cui $G$ avrebbe più bisogno di un segnale forte, ed è la fase in cui la loss originale gliene dà meno. Le due loss hanno lo stesso punto fisso ma dinamiche opposte dove conta.

\medskip
\textbf{Il collegamento con WGAN.} Il problema di fondo è che la JSD fra due distribuzioni con supporti disgiunti è costante a $\log 2$, quindi il suo gradiente è nullo: se $p_G$ e $p_{\text{data}}$ vivono su varietà che non si intersecano — situazione normale in alta dimensione — non esiste una direzione di discesa. La distanza di Wasserstein non ha questo difetto perché misura \emph{quanto} le distribuzioni sono lontane e non solo \emph{se} si sovrappongono, e fornisce gradienti utili anche a supporti disgiunti. Il critico $1$-Lipschitz (weight clipping, poi gradient penalty) serve a rendere calcolabile la forma duale di Kantorovich-Rubinstein.

\textbf{Riscontro sperimentale}: l'esercizio \texttt{12\_gan\_optimal\_d} misura questa differenza su un caso concreto e ottiene norme del gradiente di $1.4 \times 10^{-5}$ contro $1.6 \times 10^{1}$ — sei ordini di grandezza, ancora più netto del conto analitico qui sopra, perché nel caso reale $D$ è molto più sicuro di $0.01$.
\end{esempiosvolto}

\section{WGAN}

Wasserstein GAN (Arjovsky et al.\ 2017) sostituisce la JS divergence con la \emph{Earth Mover's Distance} (Wasserstein-1):
\[
W(p_{\text{data}}, p_G) = \inf_{\gamma \in \Pi(p_{\text{data}}, p_G)} \mathbb{E}_{(\vect{x},\vect{y})\sim \gamma}\|\vect{x}-\vect{y}\|.
\]
Per la dualità di Kantorovich-Rubinstein:
\[
W(p_{\text{data}}, p_G) = \sup_{\|f\|_L \leq 1} \mathbb{E}_{p_{\text{data}}}[f(\vect{x})] - \mathbb{E}_{p_G}[f(\vect{x})],
\]
con $f$ funzioni 1-Lipschitz. Il discriminatore (\emph{critic}) impara $f$:
\[
\min_G \max_{\|D\|_L \leq 1}\; \mathbb{E}_{p_{\text{data}}}[D(\vect{x})] - \mathbb{E}_{\vect{z}}[D(G(\vect{z}))].
\]
Il vincolo Lipschitz si impone via \emph{weight clipping} (originale, lento) o \emph{gradient penalty} (WGAN-GP). Vantaggio: gradiente non saturante anche quando le distribuzioni sono disgiunte.

\begin{trappola}{un GAN non si può valutare con la verosimiglianza}
Domanda tipica: ``come confronti due GAN?''. La risposta sbagliata è ``guardo la likelihood sul test set''. Un GAN è un modello \textbf{implicito}: definisce una procedura di campionamento, non una densità. Non esiste una $p_G(\vect{x})$ da valutare, né in forma chiusa né come bound — a differenza del VAE, che ha l'ELBO, e dei flow, che hanno la likelihood esatta.

Da qui nascono le metriche basate su campioni: \emph{Inception Score} e soprattutto \emph{Fréchet Inception Distance}, che confronta media e covarianza delle feature di una rete pre-addestrata sui campioni reali e generati.

\textbf{Ed è bene conoscerne i difetti, perché è lì che la domanda scava}: dipendono da una rete pre-addestrata su ImageNet, quindi non hanno senso fuori dal dominio delle immagini naturali; sono sensibili al numero di campioni; il FID assume gaussianità delle feature; e soprattutto \textbf{nessuna delle due rileva il mode collapse in modo affidabile}, perché un generatore che produce pochi modi ma di ottima qualità può ottenere un FID eccellente.

Quindi la valutazione dei GAN è un problema aperto, non una questione di scegliere la metrica giusta. È anche uno dei motivi per cui i diffusion model hanno preso il sopravvento: avendo un bound sulla verosimiglianza, offrono almeno un criterio principiato accanto alle metriche percettive.
\end{trappola}

\section{Varianti notevoli}

\begin{description}
  \item[DCGAN] convoluzioni transposte + batch norm; primo a generare immagini realistiche.
  \item[Progressive GAN] genera prima a bassa risoluzione, poi aumenta progressivamente.
  \item[CycleGAN] style transfer non appaiato: due generatori e due discriminatori con cycle consistency loss $\|F(G(\vect{x})) - \vect{x}\| + \|G(F(\vect{y})) - \vect{y}\|$.
  \item[Conditional GAN] $G(\vect{z}, \vect{y})$, $D(\vect{x}, \vect{y})$: condizionamento su label/immagine/testo.
  \item[StyleGAN] mapping network + AdaIN: controllo fine dello stile.
\end{description}

\section{Adversarial Autoencoder (AAE)}

L'AAE è il ponte VAE-GAN: l'encoder $f_\phi$ produce $\vect{z}$, e invece di forzare $\KL{q_\phi(\vect{z}|\vect{x})}{p(\vect{z})}$ \emph{in forma chiusa} (richiede $p$ tractable), si addestra un discriminatore $D$ a distinguere $\vect{z}$ campionati dal prior $p(\vect{z})$ (arbitrario!) da $\vect{z}$ prodotti da $f_\phi$:
\[
\mathcal{L}_{\text{AAE}} = \underbrace{\mathbb{E}_{q_\phi}[\log p_\theta(\vect{x}|\vect{z})]}_{\text{ricostruzione}} - \underbrace{\text{loss adversarial su } \vect{z}}_{\text{sostituisce la KL}}.
\]
Vantaggio: $p(\vect{z})$ può essere arbitrario (mistura di gaussiane, distribuzione empirica, etc.) e fornire una copertura più liscia dello spazio latente.

\begin{tcolorbox}[mybox, title={Quando usare una GAN}]
\begin{description}[leftmargin=4em,style=nextline]
  \item[Generazione di alta qualità ma rapida] sì, le GAN moderne (StyleGAN3, BigGAN) sono ancora competitive in qualità con generazione one-shot vs many-steps dei diffusion.
  \item[Inferenza, density estimation, anomaly detection via likelihood] \textbf{no}: la GAN non fornisce $p_\theta(\vect{x})$.
  \item[Conditional generation] sì (cGAN, CycleGAN, pix2pix).
  \item[Training stabile e veloce su dataset piccoli] \textbf{no}: training notoriamente instabile, richiede tuning.
  \item[Style transfer non appaiato] CycleGAN.
\end{description}
\end{tcolorbox}

\begin{tcolorbox}[mybox, title={Connessioni}]
\begin{itemize}
  \item Stesso problema del VAE (imparare $p_{\text{data}}$) con strategia opposta: implicita (sample-based) anziché esplicita (density-based).
  \item AAE $=$ VAE con KL sostituita da loss adversarial.
  \item Diffusion ha sostituito le GAN come SOTA per qualità immagine (2020+), ma le GAN restano competitive per velocità di inferenza.
  \item La distanza di Wasserstein è anche la base di flow matching e optimal transport (Cap.~\ref{ch:flowmatching}).
\end{itemize}
\end{tcolorbox}

\begin{tcolorbox}[mywarn, title={Domande d'esame — GAN}]
\begin{enumerate}
  \item Scrivi l'obiettivo minimax di Goodfellow e ricava il discriminatore ottimale $D^*$ per $G$ fisso.
  \item Mostra che, sostituendo $D^*$ nella loss, il problema è equivalente a minimizzare la divergenza di Jensen-Shannon fra $p_{\text{data}}$ e $p_G$. Qual è il punto di ottimo globale?
  \item Cos'è il \emph{mode collapse}? Cosa lo causa intuitivamente?
  \item Perché la WGAN funziona meglio della GAN classica quando le distribuzioni hanno supporti disgiunti?
  \item Qual è l'enunciato del teorema di dualità Kantorovich-Rubinstein, e come si usa per derivare la loss WGAN?
  \item In una GAN il generatore vede mai i dati reali? Spiega.
  \item Differenza concettuale fra AAE e VAE: cosa cambia nel termine di regolarizzazione e quali vantaggi/svantaggi porta.
  \item Confronta VAE e GAN su likelihood, qualità campioni, stabilità training, possibilità di inferenza.
\end{enumerate}
\end{tcolorbox}

%===============================================================
\chapter{Normalizing Flows}
\label{ch:flows}

\begin{risorsevideo}{GDL27--28 — Normalizing Flows}
\videolezione{delle lezioni GDL27 e GDL28}

\medskip
\video{https://www.youtube.com/watch?v=JBb5sSC0JoY}{CS294-158 SP20 — L3: Flow Models}{Pieter Abbeel, UC Berkeley, ~2h}\\
Cambio di variabile, coupling layer, NICE/RealNVP/Glow e flussi autoregressivi, con l'accento sul vincolo che governa tutto: rendere calcolabile il determinante dello Jacobiano.
\end{risorsevideo}

\section{Idea: cambiare variabile in modo invertibile}

Sia $\vect{z} \sim p_Z(\vect{z})$ una distribuzione base semplice (gaussiana). Cerchiamo una funzione \emph{invertibile e differenziabile} $\vect{x} = f_\theta(\vect{z})$ tale che $\vect{x} \sim p_{\text{data}}$. A differenza del VAE (che usa un decoder \emph{stocastico} e un encoder approssimato) e della GAN (che fa solo sampling), un flow \emph{computa la likelihood esatta} via cambio di variabile.

\section{Formula del cambio di variabile}

\begin{theorem}[Cambio di variabile multidimensionale]
Se $f: \mathbb{R}^d \to \mathbb{R}^d$ è $C^1$, invertibile, con Jacobiano $J_f(\vect{z}) = \partial f/\partial \vect{z}$ non singolare, allora la densità di $\vect{x} = f(\vect{z})$ è
\[
\boxed{\;p_X(\vect{x}) = p_Z(f^{-1}(\vect{x}))\, \big|\det J_{f^{-1}}(\vect{x})\big| = p_Z(\vect{z})\, \big|\det J_f(\vect{z})\big|^{-1}.\;}
\]
In log-form:
\[
\log p_X(\vect{x}) = \log p_Z(\vect{z}) - \log |\det J_f(\vect{z})|.
\]
\end{theorem}

\section{Composizione di flussi}

Sia $f = f_K \circ f_{K-1} \circ \cdots \circ f_1$, $\vect{z}_0 \sim p_Z$, $\vect{z}_k = f_k(\vect{z}_{k-1})$, $\vect{x} = \vect{z}_K$. Allora:
\[
\log p_X(\vect{x}) = \log p_Z(\vect{z}_0) - \sum_{k=1}^K \log |\det J_{f_k}(\vect{z}_{k-1})|.
\]
\textbf{Requisiti per un layer $f_k$ pratico:}
\begin{enumerate}
  \item invertibilità (per sampling: $\vect{x} \to \vect{z}$ via $f^{-1}$; per likelihood: $\vect{x} \to \vect{z}$);
  \item Jacobiano con determinante \emph{efficientemente calcolabile} (idealmente $\mathcal{O}(d)$);
  \item espressività sufficiente a comporre distribuzioni complesse.
\end{enumerate}

\begin{intuizione}{il termine di volume è un vincolo contabile}
La formula del cambio di variabile ha due pezzi e il secondo è quello che si dimentica:
\[
\log p_X(\vect{x}) = \log p_Z(\vect{z}) - \log|\det J_f(\vect{z})|.
\]

Il primo pezzo dice ``quanto era probabile il punto di partenza''. Il secondo corregge per il fatto che la trasformazione ha \emph{stirato o compresso} lo spazio in quel punto. Se $f$ espande il volume, la stessa massa di probabilità si spalma su una regione più grande e la densità cala: da qui il segno meno.

Senza quel termine la ``densità'' non integrerebbe a $1$, e il modello potrebbe barare aumentando la likelihood semplicemente comprimendo tutto verso un punto. È un vincolo di conservazione, non un dettaglio tecnico.

\textbf{Da qui discende l'intero design dei flow}: calcolare $\det J$ per una rete arbitraria costa $\mathcal{O}(d^3)$, proibitivo. Quindi tutte le architetture di flow sono risposte alla stessa domanda — \emph{come costruire trasformazioni espressive il cui Jacobiano abbia determinante calcolabile in $\mathcal{O}(d)$?} I coupling layer rispondono rendendolo triangolare; i flussi autoregressivi pure; i residual flow usano una stima stocastica; i continuous flow sostituiscono il determinante con una traccia lungo un'ODE.

\textbf{Il prezzo, che va detto all'orale}: l'invertibilità impone $\dim \vect{z} = \dim \vect{x}$. Un flow non può comprimere, quindi non ha uno spazio latente più piccolo dei dati — a differenza di VAE e latent diffusion. Ottiene in cambio la likelihood \emph{esatta}, che né il VAE (che ha un bound) né la GAN (che non ha densità) possono dare.
\end{intuizione}

\section{Layer notevoli}

\subsection{Affine flow}
$f(\vect{z}) = \mat{W}\vect{z} + \vect{b}$ con $\mat{W}$ invertibile. Determinante $|\det \mat{W}|$. Espressivo solo per casi speciali (triangolare, permutazione).

\subsection{Coupling layer (NICE, RealNVP)}
Si partiziona $\vect{z} = (\vect{z}_A, \vect{z}_B)$:
\[
\begin{aligned}
\vect{z}'_A &= \vect{z}_A,\\
\vect{z}'_B &= \vect{z}_B \odot \exp(\vect{s}_\theta(\vect{z}_A)) + \vect{t}_\theta(\vect{z}_A),
\end{aligned}
\]
con $\vect{s}, \vect{t}$ reti neurali arbitrarie (non devono essere invertibili!). Il Jacobiano è triangolare per blocchi:
\[
J = \begin{pmatrix} \mat{I} & 0 \\ * & \mathrm{diag}(\exp(\vect{s}_\theta(\vect{z}_A))) \end{pmatrix},
\quad \log|\det J| = \sum_j s_\theta(\vect{z}_A)_j.
\]
Inversione: $\vect{z}_A = \vect{z}'_A$, $\vect{z}_B = (\vect{z}'_B - \vect{t}_\theta(\vect{z}'_A)) \odot \exp(-\vect{s}_\theta(\vect{z}'_A))$. Costo $\mathcal{O}(d)$ in entrambe le direzioni.

\paragraph{NICE} ($\vect{s} \equiv 0$, solo shift): volume-preserving (det $= 1$), poco espressivo da solo ma stackabile.

\paragraph{RealNVP} affine: scale + shift. Si alterna la suddivisione (masking/squeezing) per coprire tutte le dimensioni.

\paragraph{GLOW} introduce convoluzioni $1\times 1$ invertibili $\mat{W}$ con decomposizione LU per efficienza del determinante.

\begin{esempiosvolto}{cambio di variabile in 1D e un coupling layer a due dimensioni}
\textbf{Parte 1 — verifica numerica in una dimensione.} Sia $z \sim \mathcal{N}(0,1)$ e $x = f(z) = 2z + 1$. La formula dà
\[
\log p_X(3) = \log p_Z(1) - \log|2| = \big(-\tfrac{1}{2}(1)^2 - \tfrac{1}{2}\log 2\pi\big) - \log 2 = -0.5 - 0.9189 - 0.6931 = \mathbf{-2.1121}.
\]
Verifichiamo per via diretta: $x \sim \mathcal{N}(1, 4)$, quindi
\[
p_X(3) = \frac{1}{\sqrt{8\pi}} e^{-4/8} = 0.12099, \qquad \log p_X(3) = \mathbf{-2.1121}. \quad\checkmark
\]
Il termine $-\log 2$ è esattamente la penalità per aver raddoppiato la scala: la densità si dimezza.

\medskip
\textbf{Parte 2 — perché il coupling layer è furbo.} Un coupling layer divide il vettore in due blocchi e trasforma solo il secondo, condizionando sul primo:
\[
\vect{x}_1 = \vect{z}_1, \qquad \vect{x}_2 = \vect{z}_2 \odot \exp\big(s(\vect{z}_1)\big) + t(\vect{z}_1).
\]

Con $\vect{z} = (0.5,\ -1)$, $s(0.5) = 0.4$ e $t(0.5) = 0.2$:
\[
x_1 = 0.5, \qquad x_2 = -1 \cdot e^{0.4} + 0.2 = -1.4918 + 0.2 = -1.2918.
\]

Lo Jacobiano è
\[
J = \begin{pmatrix} 1 & 0 \\ * & e^{s(\vect{z}_1)} \end{pmatrix}
\;\Rightarrow\;
\log|\det J| = s(\vect{z}_1) = 0.4.
\]

\medskip
\textbf{I tre fatti che rendono questo layer il mattone standard.} Primo: il determinante è la \emph{somma} degli $s$, calcolabile in $\mathcal{O}(d)$, perché la matrice è triangolare — e questo vale qualunque cosa siano $s$ e $t$, che possono essere reti profonde arbitrarie. Secondo: l'inversione è esplicita e altrettanto economica, $\vect{z}_2 = (\vect{x}_2 - t(\vect{x}_1)) \odot \exp(-s(\vect{x}_1))$, perché $\vect{x}_1 = \vect{z}_1$ è disponibile subito. Terzo: metà delle componenti resta ferma, quindi si alternano le partizioni fra un layer e il successivo — altrimenti quelle dimensioni non verrebbero mai trasformate.

\textbf{Nota}: $s$ e $t$ non devono essere invertibili. È il \emph{layer} a esserlo, per costruzione. È questo che permette di usare reti arbitrariamente espressive dentro un flow.
\end{esempiosvolto}

\subsection{Autoregressive flow (MAF/IAF)}
Si fattorizza $f$ in modo autoregressivo:
\[
x_i = \mu_i(\vect{x}_{1:i-1}) + z_i \exp(s_i(\vect{x}_{1:i-1})),
\]
con Jacobiano triangolare $\log|\det J| = \sum_i s_i$. \textbf{MAF} (Masked Autoregressive Flow): efficiente in \emph{density evaluation} (parallelizzabile), lento in sampling (sequenziale). \textbf{IAF} (Inverse AF): opposto. Si usano in combinazione (Parallel Wavenet) per avere entrambi efficienti.

\subsection{Residual flow}
$f(\vect{z}) = \vect{z} + g_\theta(\vect{z})$ con $g_\theta$ contraente ($\|g\|_L < 1$). Per il teorema di Banach $f$ è invertibile e l'inversa si calcola con iterazione di punto fisso $\vect{z}^{(i+1)} = \vect{z}' - g_\theta(\vect{z}^{(i)})$. Il log-determinante si stima con la serie di Neumann
\[
\log |\det J_f| = \log|\det(\mat{I} + J_g)| = \sum_{k=1}^\infty \frac{(-1)^{k+1}}{k} \mathrm{tr}(J_g^k),
\]
e $\mathrm{tr}$ si stima con \emph{Hutchinson trace estimator}: $\mathrm{tr}(\mat{A}) = \mathbb{E}_{\boldsymbol{\epsilon}\sim\mathcal{N}(0,\mat{I})}[\boldsymbol{\epsilon}^\top \mat{A} \boldsymbol{\epsilon}]$.

\subsection{Continuous Normalizing Flow (Neural ODE)}
Limite continuo di un residual flow: definendo $\vect{z}(t)$ come la soluzione di
\[
\frac{d\vect{z}(t)}{dt} = \theta_t(\vect{z}(t)),\qquad \vect{z}(0) \sim p_Z,\quad \vect{z}(T) = \vect{x},
\]
la formula di \emph{instantaneous change of variables} (Chen et al.\ NeurIPS 2018) dà
\[
\boxed{\;\log p_X(\vect{x}) = \log p_Z(\vect{z}(0)) - \int_0^T \mathrm{tr}\!\left(\frac{\partial \theta_t}{\partial \vect{z}}\right) dt.\;}
\]
\textbf{Senza determinante!} Solo la traccia del Jacobiano, e l'ODE si integra numericamente. Costo dipende dal solver (Euler, RK4, Dopri).

\begin{trappola}{likelihood esatta, e allora perché i flow non hanno vinto?}
È la domanda cattiva su questo capitolo, e ha una risposta precisa: la likelihood esatta si paga in capacità espressiva per parametro.

\emph{Vincolo di invertibilità}: ogni layer deve essere biiettivo. Questo esclude a priori tutto ciò che scarta informazione — niente ReLU non invertibili, niente pooling, niente riduzione di dimensione. Le architetture disponibili sono un sottoinsieme stretto di ciò che si userebbe altrimenti.

\emph{Vincolo di dimensione preservata}: $\dim \vect{z} = \dim \vect{x}$ ovunque. Non esiste un collo di bottiglia, quindi non esiste una rappresentazione compressa, e ogni layer lavora sempre alla piena dimensione dei dati. Su immagini ad alta risoluzione questo costa memoria e calcolo in modo diretto.

\emph{Vincolo sul determinante}: lo Jacobiano deve avere determinante calcolabile in tempo lineare, il che forza strutture triangolari (coupling, autoregressive) e quindi trasformazioni che a ogni layer toccano solo metà delle componenti. Servono molti layer per ottenere l'espressività che un VAE o un diffusion ottengono con pochi.

Il risultato pratico: a parità di budget computazionale, la qualità dei campioni resta indietro. La diffusion ha fatto la scelta opposta — rinunciare alla likelihood esatta accontentandosi di un bound — e ha guadagnato la libertà architetturale che i flow non hanno.

\textbf{La coda della storia va però raccontata}: il probability flow ODE e il flow matching recuperano l'idea del trasporto continuo \emph{senza} il vincolo di invertibilità layer per layer, perché l'invertibilità viene dall'ODE e non dall'architettura. In quel senso i flow non hanno perso, si sono trasformati.
\end{trappola}

\section{Dequantizzazione}

I dati discreti (pixel in $\{0,\dots,255\}$) violano l'assunzione di densità continua. Si aggiunge rumore uniforme $\vect{u} \in [0,1)^d$: $\vect{x}' = \vect{x} + \vect{u}$. La likelihood sul continuo è un \emph{lower bound} di quella sul discreto. Varianti più sofisticate: \emph{variational dequantization} (Flow++).

\section{Universalità}

Sotto ipotesi blande (Lipschitz, densità con supporto pieno), un flow sufficientemente profondo è un \emph{universal density approximator}: può approssimare qualunque distribuzione target con accuratezza arbitraria.

\begin{tcolorbox}[mybox, title={Quando usare un Normalizing Flow}]
\begin{description}[leftmargin=4em,style=nextline]
  \item[Density estimation esatta] sì, è l'unico modello generativo profondo che fornisce $p_\theta(\vect{x})$ esatta (no ELBO, no JS implicita).
  \item[Anomaly detection con likelihood] sì, score = $\log p_\theta(\vect{x})$ ben definito.
  \item[Generazione di alta qualità su immagini complesse] meno competitivo di diffusion/GAN moderne, ma utilizzato come testa di latent diffusion / preprocessing.
  \item[Modellare dati strutturati (audio, fisica, scienza)] sì, ad es.\ Glow per audio (Parallel Wavenet).
  \item[Conditional generation] facile (condizionare $f_\theta$).
  \item[Limite: il flow deve preservare dimensioni] non riduce, non comprime.
\end{description}
\end{tcolorbox}

\begin{tcolorbox}[mybox, title={Connessioni}]
\begin{itemize}
  \item Cambio di variabile $\leftrightarrow$ VAE: il VAE è un cambio di variabile \emph{stocastico} con dimensioni diverse; il flow è deterministico e dim-preserving.
  \item Residual flow $\to$ ODE neurale $\to$ flow matching: il limite continuo apre la strada al Cap.~\ref{ch:flowmatching}.
  \item Autoregressive flow $\leftrightarrow$ PixelRNN/PixelCNN: stessa fattorizzazione $p(\vect{x}) = \prod_i p(x_i|\vect{x}_{<i})$.
  \item Flow $\leftrightarrow$ diffusion: entrambi possono essere visti come trasporti, ma il flow è deterministico, la diffusion stocastica.
\end{itemize}
\end{tcolorbox}

\begin{tcolorbox}[mywarn, title={Domande d'esame — Normalizing Flows}]
\begin{enumerate}
  \item Scrivi e dimostra la formula del cambio di variabile multidimensionale. Spiega perché compare $|\det J|^{-1}$.
  \item Scrivi la log-likelihood di un flow composto da $K$ layer. Quali sono i requisiti pratici sui layer?
  \item Spiega come funziona un coupling layer (RealNVP). Perché il Jacobiano è triangolare? Quanto costa l'inversione?
  \item Differenza fra MAF e IAF: in quali condizioni preferiresti uno o l'altro?
  \item Cos'è un Continuous Normalizing Flow? Mostra la formula per il log-determinante (instantaneous change of variables). Vantaggi e costi.
  \item Perché serve la dequantizzazione su immagini? Cosa succede se non la fai?
  \item Confronta normalizing flow, VAE e diffusion sulla capacità di calcolare $p_\theta(\vect{x})$ esatta.
  \item Cos'è la stima Hutchinson della traccia, e perché è utile per i residual flows?
\end{enumerate}
\end{tcolorbox}

%===============================================================
\chapter{Diffusion Models}
\label{ch:diffusion}

\begin{risorsevideo}{GDL29 — Diffusion Models}
\videolezione{della lezione GDL29}

\medskip
\video{https://www.youtube.com/watch?v=DsEDMjdxOv4}{CS294-158 SP24 — L6: Diffusion Models}{Pieter Abbeel, UC Berkeley, ~2h}\\
Forward e reverse process, derivazione dell'ELBO, parametrizzazione $\epsilon$, DDIM e guidance: copre l'intero capitolo nello stesso ordine.

\medskip
\video{https://www.youtube.com/watch?v=hZ4a4NgM3u0}{Autoencoders — Deep Learning Animated}{Deepia, ~20m}\\
Da rivedere per il collegamento concettuale: il forward process è un denoising autoencoder gerarchico, e l'intuizione visiva del DAE aiuta a leggere la diffusion in quella chiave.
\end{risorsevideo}

\section{Idea: corrompere e ricostruire}

Un modello di diffusione (DDPM, Ho et al.\ 2020) definisce un processo stocastico \emph{forward} fisso che gradualmente corrompe i dati in rumore gaussiano, e impara un processo \emph{reverse} parametrico che denoisi step-by-step. Il sampling consiste nel partire da rumore puro e iterare il reverse.

\section{Forward process}

Dato $\vect{x}_0 \sim p_{\text{data}}$ e una schedule $\{\beta_t\}_{t=1}^T \subset (0,1)$:
\begin{equation}\label{eq:forward}
q(\vect{z}_t \mid \vect{z}_{t-1}) = \mathcal{N}\!\big(\sqrt{1-\beta_t}\, \vect{z}_{t-1},\, \beta_t \mat{I}\big), \qquad \vect{z}_0 = \vect{x}_0.
\end{equation}
La catena è markoviana. Definendo $\alpha_t = 1 - \beta_t$ e $\bar\alpha_t = \prod_{s=1}^t \alpha_s$, si ottiene la \emph{forma chiusa per $t$ qualsiasi}:
\begin{equation}\label{eq:diffusion-kernel}
\boxed{\;q(\vect{z}_t \mid \vect{x}_0) = \mathcal{N}\!\big(\sqrt{\bar\alpha_t}\, \vect{x}_0,\, (1-\bar\alpha_t)\mat{I}\big).\;}
\end{equation}

\begin{proof}[Derivazione]
Per induzione su $t$. Caso base $t=1$: ovvio da~\eqref{eq:forward}. Passo: se $\vect{z}_{t-1} = \sqrt{\bar\alpha_{t-1}}\vect{x}_0 + \sqrt{1-\bar\alpha_{t-1}}\boldsymbol{\epsilon}'$, allora $\vect{z}_t = \sqrt{\alpha_t}\vect{z}_{t-1} + \sqrt{\beta_t}\boldsymbol{\epsilon}'' = \sqrt{\alpha_t\bar\alpha_{t-1}}\vect{x}_0 + \sqrt{\alpha_t(1-\bar\alpha_{t-1})}\boldsymbol{\epsilon}' + \sqrt{1-\alpha_t}\boldsymbol{\epsilon}''$. Le ultime due gaussiane si combinano: $\sqrt{\alpha_t(1-\bar\alpha_{t-1}) + 1-\alpha_t}\boldsymbol{\epsilon} = \sqrt{1-\bar\alpha_t}\boldsymbol{\epsilon}$.
\end{proof}

In particolare per $T$ grande con schedule appropriata, $\bar\alpha_T \approx 0$ e $\vect{z}_T \sim \mathcal{N}(0, \mat{I})$.

\begin{intuizione}{molti problemi facili invece di uno difficile}
Il salto concettuale della diffusion sta tutto qui: generare un'immagine da zero è difficile, ma \emph{togliere un po' di rumore} da un'immagine quasi pulita è facile. Se si decompone la generazione in mille passi di denoising locale, ogni singolo passo diventa un problema di regressione banale, e la difficoltà si distribuisce lungo la catena invece di concentrarsi in un salto solo.

Da qui il resto segue in modo quasi obbligato:
\begin{itemize}[leftmargin=1.4em]
  \item il processo forward deve essere \emph{semplice e fisso}, perché non deve essere imparato: aggiungere gaussiano un po' alla volta è la scelta minima;
  \item con incrementi piccoli, la vera reverse $q(\vect{z}_{t-1}|\vect{z}_t)$ è approssimativamente gaussiana (Feller), quindi è lecito parametrizzarla come gaussiana;
  \item la distribuzione finale $q(\vect{z}_T)$ deve essere semplice e nota, altrimenti non si saprebbe da dove partire a campionare.
\end{itemize}

\textbf{Il punto sottile che l'esame cerca}: la reverse vera $q(\vect{z}_{t-1}|\vect{z}_t)$ è intrattabile, perché denoising è \emph{ambiguo} — molti dati puliti possono aver prodotto lo stesso $\vect{z}_t$, e sapere quale richiederebbe conoscere $p_{\text{data}}$. Ma se si \emph{condiziona anche su $\vect{x}_0$}, allora $q(\vect{z}_{t-1}|\vect{z}_t, \vect{x}_0)$ è gaussiana in forma chiusa. È esattamente su questo che si costruisce l'ELBO: in training $\vect{x}_0$ ce l'abbiamo, quindi possiamo confrontare due gaussiane; a inferenza non ce l'abbiamo, e per questo serve la rete.

\textbf{Il prezzo, da dire sempre}: la stessa decomposizione che rende facile ogni passo rende costosa la generazione, perché i passi vanno fatti tutti. È il trade-off strutturale della diffusion, e DDIM, distillazione e flow matching sono tre modi diversi di attaccarlo.
\end{intuizione}

\section{Reverse process}

Si parametrizza la reverse con una gaussiana:
\[
p_\theta(\vect{z}_{t-1} \mid \vect{z}_t) = \mathcal{N}\!\big(\boldsymbol{\mu}_\theta(\vect{z}_t, t),\, \sigma_t^2 \mat{I}\big).
\]
La giustificazione (Feller 1949): per $\beta_t$ piccolo, la \emph{vera} reverse $q(\vect{z}_{t-1}|\vect{z}_t)$ è approssimativamente gaussiana.

\section{Training: ELBO e parametrizzazione del rumore}

\subsection{ELBO per la catena di diffusione}
\[
\log p_\theta(\vect{x}_0) \geq \mathbb{E}_q\!\left[\log \frac{p_\theta(\vect{x}_0, \vect{z}_{1:T})}{q(\vect{z}_{1:T}|\vect{x}_0)}\right].
\]
Espandendo e usando l'identità $q(\vect{z}_t|\vect{z}_{t-1}) = q(\vect{z}_t|\vect{z}_{t-1},\vect{x}_0)$:
\[
-\text{ELBO} = \mathbb{E}_q\!\Bigg[\underbrace{\KL{q(\vect{z}_T|\vect{x}_0)}{p(\vect{z}_T)}}_{L_T \text{ (costante)}} + \sum_{t=2}^T \underbrace{\KL{q(\vect{z}_{t-1}|\vect{z}_t,\vect{x}_0)}{p_\theta(\vect{z}_{t-1}|\vect{z}_t)}}_{L_{t-1}} \underbrace{- \log p_\theta(\vect{x}_0|\vect{z}_1)}_{L_0}\Bigg].
\]
La \emph{posterior reverse condizionata su $\vect{x}_0$} è gaussiana in forma chiusa:
\[
q(\vect{z}_{t-1}|\vect{z}_t, \vect{x}_0) = \mathcal{N}(\tilde{\boldsymbol{\mu}}_t(\vect{z}_t, \vect{x}_0), \tilde\beta_t \mat{I}),
\]
con
\[
\tilde{\boldsymbol{\mu}}_t = \frac{\sqrt{\bar\alpha_{t-1}}\beta_t}{1-\bar\alpha_t} \vect{x}_0 + \frac{\sqrt{\alpha_t}(1-\bar\alpha_{t-1})}{1-\bar\alpha_t}\vect{z}_t, \qquad \tilde\beta_t = \frac{1-\bar\alpha_{t-1}}{1-\bar\alpha_t}\beta_t.
\]

\subsection{Parametrizzazione $\epsilon$}
Dalla diffusion kernel~\eqref{eq:diffusion-kernel}, $\vect{x}_0 = (\vect{z}_t - \sqrt{1-\bar\alpha_t}\boldsymbol{\epsilon})/\sqrt{\bar\alpha_t}$. Si riscrive $\tilde{\boldsymbol{\mu}}_t$ in funzione di $\boldsymbol{\epsilon}$ e si parametrizza
\[
\boldsymbol{\mu}_\theta(\vect{z}_t, t) = \frac{1}{\sqrt{\alpha_t}}\!\left(\vect{z}_t - \frac{\beta_t}{\sqrt{1-\bar\alpha_t}}\, \boldsymbol{\epsilon}_\theta(\vect{z}_t, t)\right).
\]
Il termine $L_{t-1}$ diventa allora $\propto \|\boldsymbol{\epsilon} - \boldsymbol{\epsilon}_\theta(\vect{z}_t, t)\|^2$. Ignorando i coefficienti (DDPM \emph{simplified loss}):
\begin{equation}\label{eq:ddpm-loss}
\boxed{\;\mathcal{L}_{\text{simple}}(\theta) = \mathbb{E}_{t,\vect{x}_0,\boldsymbol{\epsilon}}\!\left[\,\big\| \boldsymbol{\epsilon} - \boldsymbol{\epsilon}_\theta(\sqrt{\bar\alpha_t}\vect{x}_0 + \sqrt{1-\bar\alpha_t}\boldsymbol{\epsilon},\, t) \big\|^2 \right].\;}
\end{equation}
È una loss di \emph{denoising score matching} a $T$ livelli di rumore.

\begin{esempiosvolto}{che cosa fa davvero la schedule: $\bar\alpha_t$ e il rapporto segnale-rumore}
Prendiamo una schedule giocattolo con $T = 4$ e $\beta = (0.1,\ 0.2,\ 0.3,\ 0.4)$. Allora $\alpha_t = 1 - \beta_t$ e $\bar\alpha_t = \prod_{s \leq t}\alpha_s$:

\begin{center}
\begin{tabular}{lcccc}
\toprule
$t$ & $1$ & $2$ & $3$ & $4$ \\
\midrule
$\alpha_t$ & $0.9$ & $0.8$ & $0.7$ & $0.6$ \\
$\bar\alpha_t$ & $0.900$ & $0.720$ & $0.504$ & $0.302$ \\
$\sqrt{\bar\alpha_t}$ (segnale) & $0.949$ & $0.849$ & $0.710$ & $0.550$ \\
$\sqrt{1-\bar\alpha_t}$ (rumore) & $0.316$ & $0.529$ & $0.704$ & $0.835$ \\
\bottomrule
\end{tabular}
\end{center}

A $t = 4$ si ha $\vect{z}_4 = 0.550\,\vect{x}_0 + 0.835\,\boldsymbol{\epsilon}$ e $\mathrm{SNR}_4 = \bar\alpha_4/(1-\bar\alpha_4) = 0.434$: il rumore ha già superato il segnale.

\medskip
\textbf{Con la schedule vera.} DDPM usa $T = 1000$ e $\beta$ lineare da $10^{-4}$ a $0.02$:

\begin{center}
\begin{tabular}{lcccc}
\toprule
$t$ & $1$ & $100$ & $500$ & $1000$ \\
\midrule
$\bar\alpha_t$ & $0.9999$ & $0.897$ & $0.0786$ & $4.0 \times 10^{-5}$ \\
$\mathrm{SNR}_t$ & $10^{4}$ & $8.71$ & $0.085$ & $4.0 \times 10^{-5}$ \\
\bottomrule
\end{tabular}
\end{center}

A $t = 1000$ resta lo $0.6\%$ del segnale originale: $\vect{z}_T$ è indistinguibile da $\mathcal{N}(0,\mat{I})$, che è la condizione perché il sampling possa partire dal rumore puro.

\medskip
\textbf{Perché la schedule è un iperparametro serio.} Se il rumore cresce troppo in fretta, i primi passi di reverse devono ricostruire troppo e il problema torna difficile. Se cresce troppo lentamente, $\bar\alpha_T$ non arriva abbastanza vicino a zero e la distribuzione terminale non è davvero gaussiana, quindi il sampling parte dal punto sbagliato. L'idea guida è che i primi passi debbano \emph{corrompere i dettagli fini preservando la struttura semantica}, e gli ultimi portare tutto verso il rumore isotropo. La schedule coseno di Nichol \& Dhariwal nasce proprio dall'osservazione che quella lineare distrugge il segnale troppo presto.

\textbf{Una nota sulle varianze della reverse}: nella pratica si pone spesso $\sigma_t^2 = \beta_t$, oppure $\sigma_t^2 = \tilde\beta_t = \frac{1-\bar\alpha_{t-1}}{1-\bar\alpha_t}\beta_t$. Sono i due estremi teorici (entropia massima e minima per la reverse); alcune varianti la fanno apprendere alla rete.
\end{esempiosvolto}

\section{Sampling: DDPM e DDIM}

\subsection{DDPM (stocastico)}
\begin{align*}
\vect{z}_T &\sim \mathcal{N}(0, \mat{I}),\\
\vect{z}_{t-1} &= \frac{1}{\sqrt{\alpha_t}}\!\left(\vect{z}_t - \frac{\beta_t}{\sqrt{1-\bar\alpha_t}}\,\boldsymbol{\epsilon}_\theta(\vect{z}_t, t)\right) + \sigma_t \vect{w},\quad \vect{w} \sim \mathcal{N}(0,\mat{I}).
\end{align*}
$T$ tipicamente 1000. Lento.

\subsection{DDIM (deterministico)}
Song et al.\ 2021. Si reinterpreta la dinamica come ODE e si campiona deterministicamente con meno step (anche 50). DDIM è un caso particolare con $\sigma_t = 0$: stesso modello $\epsilon_\theta$, sampling diverso.

\section{Architettura: U-Net}

Per immagini, $\boldsymbol{\epsilon}_\theta$ è una U-Net con:
\begin{itemize}
  \item ResNet blocks down/up sampling;
  \item Self-attention nei layer di media risoluzione;
  \item \emph{Time embedding}: posizionale sinusoidale di $t$ iniettato in tutti i ResBlock;
  \item (Opzionale) condizionamento testuale via cross-attention.
\end{itemize}

\section{Conditional generation}

\subsection{Classifier guidance}
Si addestra un classificatore $p_\phi(\vect{y}|\vect{z}_t)$ \emph{su immagini rumorose}. Durante il sampling si modifica la mean:
\[
\boldsymbol{\mu}_\theta'(\vect{z}_t, t, \vect{y}) = \boldsymbol{\mu}_\theta(\vect{z}_t, t) + \gamma\, \sigma_t^2\, \nabla_{\vect{z}_t} \log p_\phi(\vect{y}|\vect{z}_t).
\]
$\gamma$ è il \emph{guidance scale}. Mappa Bayesiana: $\nabla\log p(\vect{z}_t|\vect{y}) = \nabla\log p(\vect{z}_t) + \nabla\log p(\vect{y}|\vect{z}_t)$.

\subsection{Classifier-free guidance}
Si addestra un singolo modello $\boldsymbol{\epsilon}_\theta(\vect{z}_t, t, \vect{y})$ con dropout della condizione (a volte $\vect{y} = \emptyset$). A inference:
\[
\hat{\boldsymbol{\epsilon}}(\vect{z}_t, t, \vect{y}) = (1+\gamma)\,\boldsymbol{\epsilon}_\theta(\vect{z}_t,t,\vect{y}) - \gamma\, \boldsymbol{\epsilon}_\theta(\vect{z}_t, t, \emptyset).
\]
Nessun classificatore separato. È la base di Stable Diffusion, Imagen, DALL-E 2.

\begin{remark}[Due scritture equivalenti della classifier-free guidance]
La formula usata sopra,
\[
\hat{\boldsymbol{\epsilon}} = (1+\gamma)\,\boldsymbol{\epsilon}_\theta(\vect{z}_t, t, \vect{y}) - \gamma\, \boldsymbol{\epsilon}_\theta(\vect{z}_t, t, \emptyset),
\]
si trova spesso anche nella forma a combinazione affine
\[
\hat{\boldsymbol{\epsilon}} = (1 - w)\,\boldsymbol{\epsilon}_\theta(\vect{z}_t, t, \emptyset) + w\, \boldsymbol{\epsilon}_\theta(\vect{z}_t, t, \vect{y}),
\]
che è la scrittura adottata negli handout del corso. Sono la stessa espressione con $w = 1 + \gamma$: in particolare $w = 1$ (cioè $\gamma = 0$) è il caso puramente condizionale senza guidance, e $w > 1$ amplifica la direzione condizionale. Vale la pena riconoscere entrambe le forme, perché all'orale può comparire l'una o l'altra.

\textbf{Il trade-off da nominare}: aumentando $w$ i campioni aderiscono di più al prompt ma diventano meno diversi, perché il sampling viene spinto verso i pochi modi più compatibili con la condizione. Guidance alta migliora la fedeltà e riduce la varietà — è il motivo per cui i sistemi text-to-image espongono quel parametro all'utente.
\end{remark}

\begin{remark}[Come si inietta la condizione]
Il condizionamento entra nella U-Net in modi diversi a seconda della natura del segnale: condizioni scalari o vettoriali (classe, timestep) tramite \emph{embedding sommati} dentro i blocchi residui; condizioni-immagine (inpainting, ControlNet, super-risoluzione) tramite \emph{concatenazione sui canali}; condizioni testuali tramite \emph{cross-attention} sugli embedding del testo, che è il meccanismo di Stable Diffusion e Imagen. È la stessa cross-attention della Parte~III (capitolo su encoder-decoder e cross-attention): query dalla mappa di feature dell'immagine, key e value dagli embedding testuali.
\end{remark}

\section{Latent diffusion}

Eseguire diffusion direttamente in pixel space è costoso. \emph{Latent diffusion} (Rombach et al.\ 2022, Stable Diffusion): pre-addestra un VAE che comprime $\vect{x}$ (es.\ $512\times 512$) in latente $\vect{z}$ (es.\ $64\times 64$), poi addestra diffusion in latente. Sampling: $\vect{z}_T \to \vect{z}_0 \to \vect{x} = \text{decoder}(\vect{z}_0)$.

\begin{tcolorbox}[mybox, title={Quando usare un Diffusion Model}]
\begin{description}[leftmargin=4em,style=nextline]
  \item[Generazione di immagini/audio di altissima qualità] SOTA dal 2021. Stable Diffusion, DALL-E 2, Imagen.
  \item[Conditional generation text-to-image] sì, con classifier-free guidance.
  \item[Density estimation esatta] no (ma bound ELBO).
  \item[Sampling velocissimo] \textbf{no}: 50-1000 step. Flow matching e DDIM mitigano.
  \item[Editing controllato] sì, via inpainting / SDEdit / ControlNet.
  \item[Spazio latente strutturato] meno del VAE, ma latent diffusion offre un compromesso.
\end{description}
\end{tcolorbox}

\begin{tcolorbox}[mybox, title={Connessioni}]
\begin{itemize}
  \item Forward process $=$ DAE gerarchico con $T$ livelli di rumore.
  \item Training objective $=$ denoising score matching (Cap.~\ref{ch:flowmatching}).
  \item Diffusion $=$ VAE gerarchico con encoder \emph{fisso} (non parametrizzato) e $T$ latenti della stessa dimensione di $\vect{x}$.
  \item Reverse SDE $\to$ probability flow ODE $\to$ flow matching: connessione completa in Cap.~\ref{ch:flowmatching}.
  \item Latent diffusion $=$ VAE encoder + diffusion in latente.
\end{itemize}
\end{tcolorbox}

\begin{tcolorbox}[mywarn, title={Domande d'esame — Diffusion}]
\begin{enumerate}
  \item Scrivi il forward process $q(\vect{z}_t|\vect{z}_{t-1})$ di DDPM. Deriva la forma chiusa $q(\vect{z}_t|\vect{x}_0)$ (diffusion kernel).
  \item Perché possiamo assumere che la reverse $p_\theta(\vect{z}_{t-1}|\vect{z}_t)$ sia gaussiana?
  \item Deriva l'ELBO per la catena di diffusione. Quali sono i tre tipi di termini?
  \item Mostra come si arriva alla loss simplified $\mathbb{E}\|\boldsymbol{\epsilon} - \boldsymbol{\epsilon}_\theta(\vect{z}_t, t)\|^2$ a partire dall'ELBO. Cosa si trascura?
  \item Differenza fra classifier guidance e classifier-free guidance. Vantaggi e svantaggi.
  \item Perché DDIM è più veloce di DDPM senza riaddestrare il modello?
  \item Cos'è il latent diffusion e perché è importante computazionalmente?
  \item Confronta VAE, GAN, Flow e Diffusion su: likelihood, qualità campioni, stabilità, velocità.
  \item In che senso un diffusion model è un denoising autoencoder gerarchico?
\end{enumerate}
\end{tcolorbox}

%===============================================================
\chapter{Score Matching e Flow Matching}
\label{ch:flowmatching}

\begin{risorsevideo}{GDL31 — Score matching e Flow matching}
\videolezione{della lezione GDL31}

\medskip
\video{https://www.youtube.com/watch?v=HhfLo1_yza4}{MIT 6.S184 — Lecture 03: Training Flow and Diffusion Models}{Peter Holderrieth, ~1h20m}\\
Il corso MIT 2025 costruisce flow matching e diffusion nella stessa cornice, che è esattamente il taglio di questa lezione. La lezione 3 è quella sul training, cioè su CFM e denoising score matching.

\medskip
\video{https://www.youtube.com/watch?v=DsEDMjdxOv4}{CS294-158 SP24 — L6: Diffusion Models}{Pieter Abbeel, ~2h}\\
Da rivedere per la parte SDE: la connessione fra reverse SDE e probability flow ODE è trattata con più calma che nelle slide.
\end{risorsevideo}

\section{Score function e diffusione}

Per ogni distribuzione $p(\vect{x})$ definiamo la \emph{score function}
\[
\vect{s}(\vect{x}) = \nabla_{\vect{x}} \log p(\vect{x}).
\]
È un campo vettoriale che, in ogni punto, indica la direzione di massimo incremento della log-densità. Conoscere $\vect{s}$ permette il sampling via \emph{Langevin dynamics}:
\begin{equation}
\vect{x}_{k+1} = \vect{x}_k + \frac{\eta}{2}\vect{s}(\vect{x}_k) + \sqrt{\eta}\,\boldsymbol{\epsilon}_k, \qquad \boldsymbol{\epsilon}_k \sim \mathcal{N}(0,\mat{I}),
\end{equation}
che ha $p(\vect{x})$ come distribuzione stazionaria per $\eta \to 0$.

\subsection{Vantaggio della score}

Per una densità in forma esponenziale $p(\vect{x}) = \exp(-f_\theta(\vect{x}))/Z_\theta$, la \emph{partition function} $Z_\theta$ è intrattabile in alta dimensione. Ma:
\[
\nabla_{\vect{x}} \log p(\vect{x}) = -\nabla_{\vect{x}} f_\theta(\vect{x}) - \nabla_{\vect{x}} \log Z_\theta = -\nabla_{\vect{x}} f_\theta(\vect{x}),
\]
poiché $Z_\theta$ non dipende da $\vect{x}$. La score è \emph{indipendente dalla normalizzazione}: training molto più facile per modelli energy-based.

\section{Formula di Tweedie: il ponte fra denoising e score}

Il legame fra ``predire il rumore'' e ``imparare la score'' non è una coincidenza notazionale: si ricava in tre righe dalla formula di Tweedie, ed è il primo dei due risultati di unificazione di questo capitolo.

\begin{theorem}[Tweedie]
Sia $\vect{Z}_t \mid \vect{Z}_0 = \vect{z}_0 \sim \mathcal{N}\big(\sqrt{\bar\alpha_t}\,\vect{z}_0,\ (1-\bar\alpha_t)\mat{I}\big)$ e sia $p_t$ la marginale di $\vect{Z}_t$. Allora
\[
\mathbb{E}\big[\sqrt{\bar\alpha_t}\, \vect{Z}_0 \,\big|\, \vect{Z}_t = \vect{z}_t\big] = \vect{z}_t + (1-\bar\alpha_t)\, \nabla_{\vect{z}_t} \log p_t(\vect{z}_t).
\]
\end{theorem}

Risolvendo per la stima del dato pulito:
\[
\hat{\vect{z}}_0 = \frac{\vect{z}_t + (1-\bar\alpha_t)\nabla_{\vect{z}_t}\log p_t(\vect{z}_t)}{\sqrt{\bar\alpha_t}}.
\]
Il denoising ottimo, cioè la media a posteriori del dato pulito, è governato interamente dalla score.

\begin{esempiosvolto}{da ``predici il rumore'' a ``impara la score'', in tre righe}
Dal kernel di diffusione sappiamo che $\vect{z}_t = \sqrt{\bar\alpha_t}\,\vect{z}_0 + \sqrt{1-\bar\alpha_t}\,\boldsymbol{\epsilon}_t$, da cui
\[
\vect{z}_0 = \frac{\vect{z}_t - \sqrt{1-\bar\alpha_t}\,\boldsymbol{\epsilon}_t}{\sqrt{\bar\alpha_t}}.
\]
La formula di Tweedie dà invece
\[
\vect{z}_0 = \frac{\vect{z}_t + (1-\bar\alpha_t)\nabla_{\vect{z}_t}\log p_t(\vect{z}_t)}{\sqrt{\bar\alpha_t}}.
\]
Uguagliando le due espressioni e moltiplicando per $\sqrt{\bar\alpha_t}$:
\[
\vect{z}_t - \sqrt{1-\bar\alpha_t}\,\boldsymbol{\epsilon}_t = \vect{z}_t + (1-\bar\alpha_t)\nabla_{\vect{z}_t}\log p_t(\vect{z}_t),
\]
e semplificando $\vect{z}_t$ da entrambi i lati si ottiene
\[
\boxed{\;\nabla_{\vect{z}_t} \log p_t(\vect{z}_t) = -\frac{\boldsymbol{\epsilon}_t}{\sqrt{1-\bar\alpha_t}}.\;}
\]

\medskip
\textbf{Come leggerla.} Il rumore aggiunto è proporzionale alla score cambiata di segno. Una rete addestrata a predire $\boldsymbol{\epsilon}$ sta quindi imparando, a meno di un fattore che dipende solo da $t$, il campo $\nabla \log p_t$. Le due formulazioni non sono equivalenti ``moralmente'': sono la stessa rete con una riscalatura dell'uscita.
\[
\boldsymbol{\epsilon}_\theta(\vect{z}_t, t) = -\sqrt{1-\bar\alpha_t}\; \vect{s}_\theta(\vect{z}_t, t).
\]

\textbf{Verifica del segno con un numero.} Con $\bar\alpha_t = 0.5$ si ha $\vect{s} = -\boldsymbol{\epsilon}/\sqrt{0.5} = -1.414\,\boldsymbol{\epsilon}$. Se il rumore aggiunto puntava ``verso l'alto'', la score punta verso il basso: cioè verso la regione di densità più alta, che è da dove il campione era stato allontanato. Il segno è quello giusto — la score punta sempre \emph{verso i dati}.

\textbf{Perché conta all'orale}: questa è la risposta alla domanda ``in che senso i diffusion model sono modelli score-based?''. La risposta non è ``sono simili'', è ``sono la stessa cosa, e la formula di Tweedie lo dimostra''.
\end{esempiosvolto}

\section{Score matching}

\subsection{Score matching esplicito}
\[
\mathcal{L}_{\text{SM}}(\theta) = \mathbb{E}_{\vect{x}\sim p_{\text{data}}}\!\left[\| \vect{s}_\theta(\vect{x}) - \nabla_{\vect{x}}\log p_{\text{data}}(\vect{x}) \|^2\right].
\]
Inattuabile: non conosciamo $p_{\text{data}}$.

\subsection{Implicit Score Matching (Hyvärinen 2005)}
Per integrazione per parti:
\[
\mathcal{L}_{\text{ISM}} = \mathbb{E}_{p_{\text{data}}}\!\left[ \tfrac{1}{2}\| \vect{s}_\theta(\vect{x})\|^2 + \mathrm{tr}\!\left(\nabla_{\vect{x}} \vect{s}_\theta(\vect{x})\right)\right].
\]
Costo: $\mathrm{tr}(\nabla \vect{s}_\theta)$ richiede $\mathcal{O}(d^2)$ in dimensione alta.

\subsection{Denoising Score Matching (Vincent 2011)}
Si perturba $\vect{x}$ con rumore gaussiano $\tilde{\vect{x}} = \vect{x} + \sigma \boldsymbol{\epsilon}$:
\[
\mathcal{L}_{\text{DSM}}(\theta) = \mathbb{E}_{p(\vect{x})\,p(\tilde{\vect{x}}|\vect{x})}\!\left[\| \vect{s}_\theta(\tilde{\vect{x}}) - \nabla_{\tilde{\vect{x}}}\log p(\tilde{\vect{x}}|\vect{x})\|^2\right].
\]
Per $p(\tilde{\vect{x}}|\vect{x}) = \mathcal{N}(\vect{x}, \sigma^2 \mat{I})$ si ha $\nabla_{\tilde{\vect{x}}} \log p(\tilde{\vect{x}}|\vect{x}) = -(\tilde{\vect{x}}-\vect{x})/\sigma^2 = -\boldsymbol{\epsilon}/\sigma$. Quindi
\[
\mathcal{L}_{\text{DSM}} = \mathbb{E}\!\left[\left\| \vect{s}_\theta(\tilde{\vect{x}}) + \frac{\boldsymbol{\epsilon}}{\sigma}\right\|^2\right].
\]
\textbf{È esattamente la loss DDPM riparametrizzata}: $\boldsymbol{\epsilon}_\theta(\vect{z}_t, t) = -\sqrt{1-\bar\alpha_t}\,\vect{s}_\theta(\vect{z}_t, t)$.

\section{Score-based SDE}

Song \& Ermon (NeurIPS 2019, 2021): unificano diffusion e score matching tramite SDE.
\subsection{Forward SDE}
\[
d\vect{x} = \vect{f}(\vect{x}, t)\, dt + g(t)\, d\vect{w},
\]
con $\vect{w}$ Brownian motion. Casi notevoli:
\begin{itemize}
  \item \emph{Variance Exploding} (NCSN): $\vect{f} = 0$, $g(t)$ crescente.
  \item \emph{Variance Preserving} (DDPM): $\vect{f}(\vect{x},t) = -\tfrac{1}{2}\beta(t)\vect{x}$, $g(t) = \sqrt{\beta(t)}$.
\end{itemize}

\subsection{Reverse SDE (Anderson 1982)}
Per ogni forward SDE esiste una reverse:
\begin{equation}\label{eq:reverse-sde}
d\vect{x} = \big[\vect{f}(\vect{x},t) - g^2(t) \nabla_{\vect{x}} \log p_t(\vect{x})\big] dt + g(t)\, d\bar{\vect{w}}.
\end{equation}
Il termine $\nabla \log p_t$ è la \emph{time-dependent score}, che la rete neurale apprende.

\subsection{Probability flow ODE}
Rimuovendo il termine Brownian:
\begin{equation}\label{eq:pfode}
\boxed{\;\frac{d\vect{x}}{dt} = \vect{f}(\vect{x},t) - \tfrac{1}{2}g^2(t)\nabla_{\vect{x}}\log p_t(\vect{x}).\;}
\end{equation}
Stesse \emph{marginal distributions} della SDE, ma traiettorie deterministiche $\Rightarrow$ è un Continuous Normalizing Flow! Si può:
\begin{itemize}
  \item calcolare likelihood esatta via instantaneous change of variables,
  \item invertire il flusso per encoding $\vect{x} \to \vect{z}$.
\end{itemize}

\section{Flow Matching}

L'idea (Lipman et al.\ ICLR 2023): saltare la SDE e regredire direttamente un \emph{campo di velocità} $\vect{v}_\theta(\vect{z}_\tau, \tau)$ che trasporti il prior $p_0$ nella distribuzione dei dati $p_1$ in tempo $\tau \in [0,1]$.

\begin{remark}[Convenzione di tempo]
In diffusion: $t=0$ è dato, $t=T$ è rumore. In flow matching si usa $\tau$ con $\tau=0$ rumore (sorgente), $\tau=1$ dato (target).
\end{remark}

\subsection{Probabilità path e velocità}
Si sceglie una famiglia $\{p_\tau\}_{\tau\in[0,1]}$ con $p_0$ semplice e $p_1 = p_{\text{data}}$. La \emph{marginal velocity} $\vect{u}_\tau(\vect{z})$ è il campo vettoriale che genera $p_\tau$ via l'ODE
\[
\frac{d\vect{z}_\tau}{d\tau} = \vect{u}_\tau(\vect{z}_\tau).
\]
L'obiettivo ideale è
\[
\mathcal{L}_{\text{FM}} = \mathbb{E}_{\tau, \vect{z}_\tau \sim p_\tau}\!\left[ \|\vect{v}_\theta(\vect{z}_\tau, \tau) - \vect{u}_\tau(\vect{z}_\tau)\|^2\right],
\]
ma $\vect{u}_\tau$ non è osservabile.

\subsection{Conditional Flow Matching (CFM)}
Si sceglie un \emph{conditional path} $q_\tau(\vect{z}|\vect{x})$ con velocità $\vect{u}_\tau(\vect{z}|\vect{x})$ noti. Loss:
\begin{equation}\label{eq:cfm}
\boxed{\;\mathcal{L}_{\text{CFM}}(\theta) = \mathbb{E}_{\tau,\vect{x},\vect{z}\sim q_\tau(\cdot|\vect{x})}\!\left[\|\vect{v}_\theta(\vect{z},\tau) - \vect{u}_\tau(\vect{z}|\vect{x})\|^2\right].\;}
\end{equation}

\begin{theorem}[Lipman et al.\ 2023]
Sotto ipotesi tecniche, $\nabla_\theta \mathcal{L}_{\text{CFM}} = \nabla_\theta \mathcal{L}_{\text{FM}}$. Le velocità condizionate \emph{si mediano} sulla velocità marginale corretta.
\end{theorem}
\begin{proof}[Sketch]
Si mostra che $\vect{u}_\tau(\vect{z}) = \mathbb{E}[\vect{u}_\tau(\vect{z}|\vect{x}) \mid \vect{z}_\tau = \vect{z}]$. La differenza fra le due loss è una costante indipendente da $\theta$.
\end{proof}

\subsection{Path notevoli}

\paragraph{Rectified flow.}
\[
\vect{z}_\tau = (1-\tau)\vect{z}_0 + \tau\, \vect{x}, \quad \vect{z}_0 \sim p_0, \quad \vect{x} \sim p_{\text{data}}.
\]
Velocità conditional costante:
\[
\vect{u}_\tau(\vect{z}_\tau | \vect{x}, \vect{z}_0) = \frac{d\vect{z}_\tau}{d\tau} = \vect{x} - \vect{z}_0.
\]
Traiettorie rettilinee. Sampling: pochi step di Euler bastano.

\paragraph{Gaussian probability path.}
\[
\vect{z}_\tau = \alpha_\tau\vect{x} + \sigma_\tau\boldsymbol{\epsilon}, \qquad \alpha_0=0,\,\sigma_0=1,\, \alpha_1=1,\,\sigma_1=0.
\]
\[
\vect{u}_\tau(\vect{z}_\tau|\vect{x},\boldsymbol{\epsilon}) = \dot\alpha_\tau \vect{x} + \dot\sigma_\tau \boldsymbol{\epsilon}.
\]
Per appropriate scelte di $\alpha_\tau, \sigma_\tau$ si recuperano le path della diffusion. \emph{Flow matching gaussiano e diffusion sono lo stesso modello sotto parametrizzazioni e pesi diversi.}

\section{Sampling}

Integrazione dell'ODE:
\[
\frac{d\vect{z}_\tau}{d\tau} = \vect{v}_\theta(\vect{z}_\tau, \tau), \qquad \vect{z}_0 \sim p_0.
\]
Solver Euler:
\[
\vect{z}_{\tau+\Delta\tau} = \vect{z}_\tau + \Delta\tau\, \vect{v}_\theta(\vect{z}_\tau, \tau).
\]
Numero di step tipicamente molto inferiore a DDPM (anche 4-8 step con rectified flow + distillation).

\begin{esempiosvolto}{rectified flow: la velocità è costante, e si vede}
Caso unidimensionale. Sorgente $z_0 = 0.5$, dato $x = 2.0$, cammino lineare
\[
z_\tau = (1-\tau)z_0 + \tau x.
\]

\textbf{Il bersaglio della regressione.}
\[
u_\tau = \frac{dz_\tau}{d\tau} = x - z_0 = 2.0 - 0.5 = 1.5,
\]
costante rispetto a $\tau$. Non serve né una score, né un kernel reverse, né un'ELBO: la coppia (sorgente, dato) determina direttamente il bersaglio.

\medskip
\textbf{Sampling con Euler, $K = 4$ passi, $\Delta\tau = 0.25$.} Supponendo che la rete abbia imparato $v_\theta \equiv 1.5$:
\[
0.5 \;\to\; 0.875 \;\to\; 1.25 \;\to\; 1.625 \;\to\; 2.0.
\]
Quattro passi, arrivo \emph{esatto}. Non è una coincidenza: il metodo di Euler è esatto quando la traiettoria è una retta, perché l'errore di troncamento dipende dalla derivata seconda, che qui è nulla.

\medskip
\textbf{Ecco perché flow matching campiona in pochi passi.} Un DDPM ha bisogno di centinaia o migliaia di passi perché le sue traiettorie sono curve e stocastiche, e ogni passo di Euler introduce un errore. Rectified flow costruisce deliberatamente traiettorie il più possibile rettilinee, e su traiettorie rettilinee l'integratore non sbaglia. Le procedure di \emph{reflow} iterano la costruzione per raddrizzare ulteriormente i cammini, arrivando a generare in $1$--$4$ passi.

\textbf{L'onestà del caso reale}: con un solo punto le traiettorie sono banalmente rette. Su una distribuzione vera il campo marginale $\vect{u}_\tau(\vect{z})$ è la \emph{media} delle velocità condizionali che passano per $\vect{z}$, e quella media in generale curva. Il teorema del CFM garantisce che regredire i bersagli condizionali dia il gradiente giusto per il campo marginale — non che il campo marginale sia rettilineo.
\end{esempiosvolto}

\begin{tcolorbox}[mybox, title={Le due ricette di training, affiancate}]
La distanza concettuale fra diffusion e flow matching è grande; quella fra le due implementazioni è minima. Vale la pena saperlo dire, perché è una domanda che si presta bene a distinguere chi ha capito da chi ha memorizzato.

\medskip
\begin{tabularx}{\linewidth}{@{}p{0.46\linewidth}p{0.46\linewidth}@{}}
\textbf{Diffusion} & \textbf{Flow matching} \\[2pt]
\hline\\[-6pt]
1. campiona $\vect{x} \sim p_{\text{data}}$ & 1. campiona $\vect{x} \sim p_{\text{data}}$ \\
2. campiona il tempo $t$ & 2. campiona la sorgente $\vect{z}_0 \sim p_0$ (o $\boldsymbol{\epsilon}$) \\
3. campiona $\boldsymbol{\epsilon} \sim \mathcal{N}(0,\mat{I})$ & 3. campiona il tempo $\tau$ \\
4. costruisci $\vect{z}_t = \sqrt{\bar\alpha_t}\vect{x} + \sqrt{1-\bar\alpha_t}\boldsymbol{\epsilon}$ & 4. costruisci $\vect{z}_\tau = (1-\tau)\vect{z}_0 + \tau\vect{x}$ \\
5. predici $\boldsymbol{\epsilon}$ (o la score) & 5. predici la velocità $\vect{v}$ \\
\end{tabularx}

\medskip
Stessa struttura, stessa backbone (U-Net o Transformer con time embedding), stessa loss quadratica. Cambia \emph{il bersaglio della regressione}: $\boldsymbol{\epsilon}$, $\vect{s}$ oppure $\vect{v}$. E cambia chi decide il cammino: nella diffusion lo detta il processo di rumore, nel flow matching lo si sceglie.
\end{tcolorbox}

\section{Riepilogo unificante}

\begin{tcolorbox}[mybox, title={Quattro lenti sullo stesso oggetto}]
Lo stesso modello generativo si può vedere come:
\begin{itemize}
  \item \textbf{Diffusion} (DDPM): denoising step-by-step di una catena di Markov.
  \item \textbf{Score matching}: regressione del campo $\nabla\log p_t$.
  \item \textbf{Probability flow ODE}: continuous normalizing flow deterministico.
  \item \textbf{Flow matching}: regressione del campo di velocità lungo un cammino di probabilità.
\end{itemize}
La differenza pratica è nella \emph{parametrizzazione della testa} (predice $\boldsymbol{\epsilon}$, $\vect{s}$ o $\vect{v}$) e nelle scelte di path/scheduling, non nell'architettura.
\end{tcolorbox}

\begin{tcolorbox}[mybox, title={Quando usare Flow Matching}]
\begin{description}[leftmargin=4em,style=nextline]
  \item[Generazione veloce ad alta qualità] sì, $\mathcal{O}(10)$ step vs $\mathcal{O}(1000)$ di DDPM.
  \item[Path flessibili per dati con struttura nota] sì, si può scegliere il path che fitta il problema (rectified, OT, schroedinger bridge).
  \item[Optimal transport] connessione naturale, rectified flow approssima OT.
  \item[Density estimation esatta] sì, tramite probability flow ODE.
  \item[Quando preferire diffusion classico?] solo per maggiore maturità del software / pretrained models. Concettualmente flow matching è preferibile.
\end{description}
\end{tcolorbox}

\begin{tcolorbox}[mywarn, title={Domande d'esame — Score \& Flow Matching}]
\begin{enumerate}
  \item Cos'è la score function di una distribuzione e perché è utile per modelli energy-based?
  \item Deriva la loss del Denoising Score Matching per rumore gaussiano. Mostra che è equivalente alla loss simplified di DDPM.
  \item Scrivi forward e reverse SDE. Cosa rappresenta il termine $g^2(t)\nabla\log p_t$ nella reverse?
  \item Cos'è la probability flow ODE e perché collega diffusion e normalizing flows?
  \item Spiega l'idea di Flow Matching: cosa si regredisce e perché si usa un \emph{conditional} flow matching?
  \item Mostra il path rectified flow e calcola la velocità conditional. Perché le traiettorie sono rettilinee?
  \item In che senso flow matching e diffusion sono lo stesso modello sotto diverse parametrizzazioni?
  \item Vantaggi di flow matching in termini di numero di passi di sampling.
\end{enumerate}
\end{tcolorbox}

%===============================================================
\chapter{Causal Representation Learning}
\label{ch:crl}

\begin{risorsevideo}{GDL30 — Disentanglement e Causal Representation Learning}
\videolezione{della lezione GDL30}

\medskip
\video{https://www.youtube.com/watch?v=f8JrbaTR1vg}{UAI 2023 Tutorial: Causal Representation Learning}{~2h}\\
Copre il risultato di impossibilità di Locatello, l'identificabilità tramite interventi e i modelli temporali tipo CITRIS: è la mappa completa del capitolo, tenuta da chi lavora sul tema.

\medskip
\video{https://www.youtube.com/watch?v=AuZu0L0PEgk}{Causal Models}{Brady Neal, ~1h}\\
Da rivedere per riallacciare la Parte~I: il CRL è la causalità del Cap.~\ref{ch:causality} applicata a variabili latenti anziché osservate.
\end{risorsevideo}

\section{Disentanglement: definizione e motivazione}

Si assume che le osservazioni $\vect{x}$ siano generate da \emph{fattori di variazione} $\vect{s} = (s_1, \dots, s_k)$ \emph{indipendenti}:
\[
\vect{x} = f^*(\vect{s}), \qquad p(\vect{s}) = \prod_j p(s_j).
\]
Una rappresentazione $\vect{z} = h(\vect{x})$ è \emph{disentangled} se ciascuna componente $z_j$ corrisponde a un unico fattore $s_{\pi(j)}$ (permutazione $\pi$) e ad un'eventuale trasformazione monotona componente-per-componente. Intervenire su $z_j$ cambia solo $s_{\pi(j)}$.

Disentanglement è auspicabile per: interpretabilità, equità, generalizzazione compositiva, transfer learning.

\section{Approcci unsupervised}

\subsection{$\beta$-VAE (Higgins et al.\ 2017)}
\[
\mathcal{L}_\beta = \mathbb{E}_{q_\phi}[\log p_\theta(\vect{x}|\vect{z})] - \beta\,\KL{q_\phi(\vect{z}|\vect{x})}{p(\vect{z})},\qquad \beta > 1.
\]
La regolarizzazione forte spinge $q_\phi(\vect{z}|\vect{x})$ verso il prior fattorizzato $p(\vect{z}) = \mathcal{N}(0,\mat{I})$, sperando in componenti indipendenti.

\subsection{FactorVAE (Kim \& Mnih 2018)}
Decompone la KL e penalizza specificamente la \emph{total correlation} (TC) del marginale $q_\phi(\vect{z}) = \mathbb{E}_{\vect{x}}[q_\phi(\vect{z}|\vect{x})]$:
\[
\text{TC}(\vect{z}) = \KL{q_\phi(\vect{z})}{\prod_j q_\phi(z_j)}.
\]
La TC si stima con un discriminatore.

\subsection{$\beta$-TCVAE (Chen et al.\ 2018)}
Decompone $\KL{q_\phi(\vect{z}|\vect{x})}{p(\vect{z})}$ in tre termini (mutual info $\vect{x}\leftrightarrow\vect{z}$, TC, KL dimensionwise) e pesa solo quello che penalizza la TC.

\begin{intuizione}{perché il disentanglement non supervisionato è impossibile}
L'argomento di Locatello è più semplice di quanto la sua fama suggerisca, e conviene saperlo raccontare in tre passaggi.

Supponiamo di avere un modello generativo con latente $\vect{z} \sim \mathcal{N}(0, \mat{I})$ e decoder $g$, che riproduce perfettamente la distribuzione dei dati. Prendiamo ora una rotazione $\mat{R}$ dello spazio latente e definiamo $\tilde{\vect{z}} = \mat{R}\vect{z}$, $\tilde g = g \circ \mat{R}^{-1}$.

Poiché la gaussiana standard isotropa è invariante per rotazioni, $\tilde{\vect{z}}$ ha \emph{esattamente} la stessa distribuzione di $\vect{z}$, e il nuovo modello genera \emph{esattamente} la stessa distribuzione sui dati. I due modelli sono indistinguibili da qualunque quantità calcolata sui dati: stessa likelihood, stesso ELBO, stessa qualsiasi cosa.

Ma se $\vect{z}$ era disentangled — una coordinata per il colore, una per la posizione — allora $\tilde{\vect{z}}$ mescola colore e posizione in ogni coordinata. Le due rappresentazioni hanno lo stesso valore per ogni criterio osservativo e opposto valore di disentanglement. Quindi \emph{nessun} obiettivo non supervisionato può preferire l'una all'altra: le rotazioni sono un'intera famiglia infinita di soluzioni equivalenti.

\textbf{Da cui la morale, che è costruttiva e non nichilista}: serve rompere quella simmetria dall'esterno. Le tre vie sono supervisione debole (coppie, gruppi, etichette parziali), struttura temporale (TCL: se le sorgenti evolvono in modo indipendente nel tempo, la storia le identifica), oppure \emph{interventi} — ed è esattamente qui che la causalità entra nel discorso. Un intervento su un fattore cambia una sola coordinata vera e lascia ferme le altre; sotto rotazione ne cambierebbe molte. L'intervento distingue ciò che l'osservazione non distingue, esattamente come nella Parte~I.

\textbf{Nota su $\beta$-VAE}: alza $\beta$ e ottiene rappresentazioni più fattorizzate, ma questo non contraddice il teorema. Sta introducendo un bias induttivo implicito (verso assi allineati alle coordinate), non identificando i fattori veri — e infatti il risultato dipende fortemente dal seed, che è proprio la verifica sperimentale di Locatello.
\end{intuizione}

\section{Impossibilità (Locatello et al.\ 2019)}

\begin{theorem}[Indeterminatezza non-supervisionata]
Per qualunque $\vect{z}$ con componenti marginalmente indipendenti, esiste un'infinità di trasformazioni $\vect{z}' = f(\vect{z})$ che preservano la indipendenza marginale ma producono rappresentazioni \emph{entangled} rispetto ai fattori ground-truth, con \emph{stessa loss ELBO}.
\end{theorem}

In altre parole, marginale indipendenza è necessaria ma non sufficiente. Il successo empirico di $\beta$-VAE e FactorVAE dipende da \emph{inductive bias} (architettura, prior, regolarizzazione), non da garanzie formali.

\section{Disentanglement supervisionato debolmente}

Per ottenere \emph{identifiability} servono informazioni ausiliarie $\vect{u}$ (tempo, ambiente, classe) tali che $p(\vect{s}|\vect{u})$ vari con $\vect{u}$.

\subsection{Time-Contrastive Learning (TCL, Hyvärinen \& Morioka 2016)}
Si divide una serie temporale in segmenti e si allena un classificatore a riconoscere il segmento da $\vect{x}$. Sotto ipotesi (sorgenti exp-family, mixing iniettivo, segmenti "abbastanza diversi"), si recuperano le sorgenti a meno di permutazione e trasformazione monotona componente-per-componente.

\subsection{iVAE (Khemakhem et al.\ 2020)}
Variante condizionale del VAE in cui il prior dipende dalla variabile ausiliaria:
\[
p_\theta(\vect{z}|\vect{u}) = \prod_j p_j(z_j | \boldsymbol{\lambda}(\vect{u})).
\]
ELBO con $p(\vect{z}|\vect{u})$ al posto di $p(\vect{z})$:
\[
\mathcal{L}_{\text{iVAE}} = \mathbb{E}_{q_\phi}[\log p_\theta(\vect{x}|\vect{z})] - \KL{q_\phi(\vect{z}|\vect{x},\vect{u})}{p_\theta(\vect{z}|\vect{u})}.
\]
Identifiability: sotto condizioni su $\boldsymbol{\lambda}$ e sull'iniettività del decoder, si recuperano le sorgenti a meno di trasformazioni componente-wise.

\section{Causal Representation Learning}

\subsection{Setup}
I fattori $\vect{s}$ non sono indipendenti ma seguono uno \emph{Structural Causal Model} (SCM) (Pearl, Cap.~\ref{ch:causality}):
\[
s_i := f_i(\mathrm{Pa}(s_i), \epsilon_i), \quad \epsilon_i \perp \epsilon_j.
\]
Si osserva $\vect{x} = h^*(\vect{s})$ con $h^*$ mixing iniettivo. L'obiettivo è recuperare sia le variabili causali $\vect{s}$ sia il loro \emph{causal graph} $G$ a partire dalle osservazioni.

\subsection{Identifiability via interventi}
Senza informazione extra il problema è impossibile (analogamente alla causalità classica): le \emph{Markov equivalence classes} rendono indistinguibili più grafi. Servono \emph{interventi}.

\subsection{CITRIS (Lippe et al.\ 2022)}
Setting: \emph{Temporal Causal Triplets} $(\vect{x}_t, \vect{x}_{t+1}, \vect{I}_{t+1})$ dove $\vect{I}_{t+1}$ è una maschera binaria che indica quali variabili causali sono state intervenute al tempo $t+1$ (non si conoscono i valori dell'intervento).

Architettura:
\begin{itemize}
  \item Encoder $\vect{z}_t = g_\phi(\vect{x}_t)$;
  \item Assignment function $\boldsymbol{\psi}: \mathcal{Z} \to \{1,\dots,k\}^m$ che mappa ogni dimensione latente a una variabile causale;
  \item Transition prior $p(\vect{z}_{t+1} | \vect{z}_t, \vect{I}_{t+1})$ parametrizzato in modo che ogni gruppo $\vect{z}^{(i)}$ dipenda da $\vect{z}_t$ tramite il proprio modulo, modulato da $I_i$;
  \item Decoder $p(\vect{x}_t | \vect{z}_t)$.
\end{itemize}
Training: ELBO temporale + (opzionale) target classifier che predice $\vect{I}_{t+1}$ dai latenti.

\begin{theorem}[Identifiability CITRIS]
Nel limite di dati infiniti, CITRIS recupera le variabili causali ground-truth se per ogni $C_i$:
\begin{enumerate}
  \item esiste un regime in cui $C_i$ è intervenuto e non sempre congiunto a un altro $C_j$;
  \item $C_i$ dipende sempre dalla sua variabile di intervento $I_i$.
\end{enumerate}
\end{theorem}

\subsection{Varianti}
\begin{description}
  \item[CITRIS-NF] disentangle una rappresentazione pre-addestrata via normalizing flows.
  \item[BISCUIT] generalizza a interventi binari sconosciuti.
\end{description}

\section{Triplet evaluation}

Test: dato un modello CITRIS addestrato, si genera un \emph{interchange intervention} sostituendo il valore di una variabile causale da un'immagine con quello dell'altra. Se la generazione cambia in modo coerente solo sulla variabile sostituita, il modello è disentangled.

\begin{tcolorbox}[mybox, title={Quando usare CRL}]
\begin{description}[leftmargin=4em,style=nextline]
  \item[Sistemi con interventi disponibili] (RL, scienze sperimentali) sì.
  \item[Dati osservazionali puri] CRL puro è impossibile (Locatello). Si può ancora usare $\beta$-VAE per interpretabilità empirica, senza garanzie.
  \item[Time-series con regimi non stazionari] iVAE/TCL.
  \item[Generative model che generalizzi a interventi distribuzionali] sì, CRL è il giusto framework.
\end{description}
\end{tcolorbox}

\begin{tcolorbox}[mybox, title={Connessioni}]
\begin{itemize}
  \item $\beta$-VAE $\to$ FactorVAE $\to$ $\beta$-TCVAE: progressione di obiettivi per disentanglement.
  \item Locatello 2019 $\to$ TCL/iVAE $\to$ CITRIS: dalla impossibilità unsupervised al requirements di interventi.
  \item CITRIS $=$ VAE temporale con prior strutturato per intervento (parente del CVAE, Cap.~\ref{ch:vae}).
  \item Connette Causalità (Cap.~\ref{ch:causality}) e Deep Generative Learning.
\end{itemize}
\end{tcolorbox}

\begin{tcolorbox}[mywarn, title={Domande d'esame — Causal Representation Learning}]
\begin{enumerate}
  \item Definisci formalmente cosa significa che una rappresentazione è \emph{disentangled}.
  \item Cos'è il $\beta$-VAE e perché aumentare $\beta$ dovrebbe promuovere il disentanglement?
  \item Enuncia il risultato di impossibilità di Locatello et al.\ 2019. Quale assunto viene violato dall'esistenza di trasformazioni che preservano la indipendenza marginale?
  \item Cos'è una variabile ausiliaria $\vect{u}$ in iVAE/TCL? Perché permette di superare l'impossibilità?
  \item Cos'è il problema di Causal Representation Learning? In cosa differisce da disentanglement classico?
  \item Descrivi il setup di CITRIS (input, output, ipotesi su interventi).
  \item Cosa garantisce il teorema di identifiability di CITRIS?
  \item Differenza pratica fra CITRIS-VAE e CITRIS-NF.
  \item Confronta causalità classica (Cap.~\ref{ch:causality}) e CRL: quale è il contributo del deep learning?
\end{enumerate}
\end{tcolorbox}

%===============================================================
\chapter*{Confronto finale Parte IV}
\addcontentsline{toc}{chapter}{Confronto finale Parte IV}

\begin{tcolorbox}[mybox, title={Confronto modelli generativi profondi}, breakable]
\small
\begin{tabular}{p{2.6cm} p{1.5cm} p{1.5cm} p{1.6cm} p{1.7cm} p{1.7cm}}
\toprule
\textbf{Proprietà} & \textbf{VAE} & \textbf{GAN} & \textbf{Flows} & \textbf{Diffusion} & \textbf{Flow Match.}\\
\midrule
Likelihood esatta & ELBO & no & sì & ELBO/no & no\\
Score implicita & no & no & sì (PF-ODE) & sì & sì\\
Qualità campioni & media & alta & media & molto alta & alta\\
Velocità sampling & veloce & veloce & veloce & lenta ($T\!\sim\!10^3$) & media ($T\!\sim\!10$)\\
Stabilità training & stabile & instabile & stabile & stabile & stabile\\
Latente strutturato & sì & no & sì & parziale & sì\\
Cond.\ generation & facile & possibile & facile & facile (CFG) & facile\\
Encoder esplicito & sì & no & sì (inverso) & no/fisso & no\\
Dim.\ latente & $\neq$ dato & $\neq$ dato & $=$ dato & $=$ dato & $=$ dato\\
Anno (paper seminale) & 2013 & 2014 & 2014--15 & 2020 & 2023\\
\bottomrule
\end{tabular}
\end{tcolorbox}

\bigskip

\begin{tcolorbox}[mybox, title={Famiglia generativa — la storia in 5 passi}]
\begin{enumerate}
  \item \textbf{Autoencoder} (1980s, deep AE 2006): latente, ma non probabilistico.
  \item \textbf{VAE} (Kingma 2013): aggiunge un encoder stocastico, ELBO, reparametrization $\Rightarrow$ il primo modello latente differenziabile end-to-end.
  \item \textbf{GAN} (Goodfellow 2014): cambia paradigma — impara a campionare, non a stimare la densità.
  \item \textbf{Normalizing Flow} (Dinh 2014, 2017): cambio di variabile esplicito, likelihood esatta, ma vincoli di invertibilità.
  \item \textbf{Diffusion} (Ho 2020): VAE gerarchico con encoder fisso a $T$ livelli, training stabile, qualità SOTA, ma sampling lento.
  \item \textbf{Flow Matching} (Lipman 2023): regressione diretta del campo di velocità $\Rightarrow$ diffusion senza SDE, sampling efficiente.
\end{enumerate}
Le 5 lenti unificanti del corso si manifestano qui:
\begin{itemize}
  \item \emph{Message passing}: nessuno (questi modelli operano su tensor euclidei).
  \item \emph{Generative family} (questo capitolo).
  \item \emph{Attention}: usato come building block (U-Net dei diffusion, Transformer per latent diffusion).
  \item \emph{EM/Latent variable}: il VAE è il discendente diretto di EM (Cap.~\ref{ch:em}).
  \item \emph{Causality}: CRL applica il framework generativo al recupero di SCM.
\end{itemize}
\end{tcolorbox}
